\section{The seesaw identity}\label{sec:seesaw}

This section proves that the lower chain of the Poincar\'e series of a ramified source form is a fixed
nonzero multiple of a composite lift of the form itself (\cref{thm:seesaw}). The mechanism is Kudla's seesaw
\cite{Kudla1984}: one pairs $\iota_\vartheta(\mathcal P\varphi)$ with a form on $Y_m$, moves the $m$ lifts to
the other side as theta descents, and unfolds the Bessel period of $\varphi$ against them one lift at a time
(\cref{ss:B-unfolding}). Each unfolding step introduces an \emph{outer variable}, a point of $\U(Y_j)(\Ad)$;
its archimedean part is integrated once and for all, and its finite part runs over a finite set of classes,
the \emph{cut}, with weights. Two consecutive unfolding steps combine into one double lift between the
definite parts $X_{2i}\to X_{2i+1}\to X_{2i+2}$, and the finite part of the datum of a double lift, once
evaluated, is a combination of the glued data of \cref{ss:composite-lift}. Three points need control, and
each has a subsection: the support of the unfolded functions (\cref{ss:support}), which makes the cuts
finite with multiplicities dividing powers of the norm $q$ of the modulus of $\vartheta$; the evaluation of
the vacuum along a double lift (\cref{ss:neighbour}), which produces the gluing above $N$; and the
denominators and sizes of all coefficients (\cref{ss:admissible}), which are those of a chain datum.

Throughout, $\vartheta$ is a tame character with modulus $Q$, $q=\prod_{w\in Q}\Nm w$, and $\mathfrak q$ is a
finite idele generating the ideal $\prod_{w\in Q}w$ (the generator of $\pi_MQ$ divided by that of $\pi_M$). The
levels $J_j=J_\vartheta(Y_j)$ are the elements of the lattice stabiliser $J^\circ_j$ of $Y_j$ that reduce to $1$ at
the level places, and $J_M\subseteq J^\circ_0$ is the chosen level of $Y_0$.
\lean{New.Data.twistIdele, New.Data.valued_twistIdele, New.Data.notMem_Q_of_mem_piNPlaces, New.Data.chosenLevel_inf_trivialOn_le, New.Data.lt_m_of_le}

\subsection{Graded spaces and their lattices}\label{ss:graded}
A \emph{grading} of a hermitian space $V$ with a point is a splitting of its basis into a part $X$, on which the
Gram matrix is the identity, and a \emph{string} $b_0,\dots,b_{s-1}$ orthogonal to $X$, with $\langle b_i,b_{s-1-i}\rangle=1$
for $2i\ne s-1$ and $\langle b_i,b_i\rangle=-N$ for the middle vector when $s$ is odd. The \emph{grade} is $0$ on $X$ and
$s-1-2i$ on $b_i$; a matrix \emph{raises grades} if its entry $(p,p')$ vanishes unless the grade of $p$ exceeds
that of $p'$. The \emph{unipotent group} of the grading consists of the points whose matrix minus $1$ raises
grades, and the \emph{Bessel group} $R_V=U_V\rtimes\U(X)$ is generated by it and by the embedded unitary group of
$X$; it is the group of \cref{ss:B-bessel}, and taking the $X$-block is a homomorphism onto $\U(X)$ with kernel
$U_V$. For a grading $Y'$ of $V'$ with a string of length $s-1$, the \emph{shift} $\mathrm{sh}\colon V\to V'$ kills $X$
and sends $b_i$ to a fixed multiple of $b'_i$ (to $0$ for $i=s-1$); it lowers grades by one, and $\mathrm{sh}^\dagger\mathrm{sh}$
and $\mathrm{sh}\,\mathrm{sh}^\dagger$ are the shifts of $V$ and $V'$ when the constants of the two gradings agree.
\lean{New.Grading, New.Grading.stringMap, New.Grading.besselGroup, New.Grading.unipotent, New.Grading.inclX, New.Grading.stringMap_eq_zero_of_gradeOf_ne, New.Grading.stringProj_stringMapFin, New.Grading.stringMapFin_stringProj, New.Seesaw.inclX_mul_stringMap_mul_inclX_inv, New.Seesaw.formAdjoint_fromRows}

For finite ideles $a$, $\pi$ and a norm-one adele $u$, the \emph{graded lattice basis} is obtained from the basis
of $V$ as follows. If $s$ is odd, multiply the string vectors by $u$ and replace the last vector $e_0$ of
$X\oplus\langle e_0\rangle$ by $e_0/\bar\pi+uz/(\pi\bar\pi)$, $z$ the middle vector; if $s$ is even, multiply the first half of the
string by $u/a$ and the second half by $u\bar a$. The point of $V$ \emph{spans the graded lattice} if its frame and
this basis span the same $\hat\OO_E$-lattice $\Lambda$. For a finite point $x$ write $\Lambda_x=x\Lambda$, $\Lambda^+_x$ for its
projection to $X^\perp$ and $\Lambda^-_x=\Lambda_x\cap X^\perp$, and drop the index when $x=1$. A lattice is \emph{graded} if it
contains the part in each grade of each of its vectors.
\lean{New.gradedLatticeBasis, New.Grading.SpansLattice, New.Seesaw.gluedGram, New.Seesaw.map_gluedGram, New.Seesaw.gluedGram_mul_eq_one, New.Seesaw.map_gradedGram, New.Seesaw.isHermitian_map_gradedGram_natCast, New.Seesaw.eq_rev_iff_of_odd, New.Seesaw.ne_rev_of_not_odd, New.eq_rev_iff, New.eq_of_eq_rev, New.odd_of_eq_rev, New.lineIdele, New.powerIdele, New.powerIdele_eq_mul, New.powerIdele_mul_zpow, New.lineIdele_divStar_pow, New.toUnits_divStar_pow}

\begin{proposition}[The graded lattice]\label{prop:graded-lattice}
Let $N=\pi\bar\pi$, with $\pi$ a uniformiser at the places of a set $P$ containing no conjugate pair and a unit
elsewhere.
\begin{enumerate}[label=\textup{(\alph*)}]
\item The Gram matrix of the graded lattice basis is hermitian and invertible over every subring; over $\C$ the
form is congruent to a diagonal of signs, $\lfloor(s+1)/2\rfloor$ of them negative.
\item If the point spans the graded lattice, $\Lambda$ is graded, and $\Lambda^-$, $\Lambda^+$ are sums of lines on the string
vectors whose generators differ only by $N^{-1}$ at the middle vector. A vector of $\Lambda^+$ outside $\Lambda^-$ exists
only for odd $s$, and it generates $\Lambda^+$ over $\Lambda^-$ at some place above $N$.
\item If the constants of $Y$ and $Y'$ have the valuations of $a\bar a$ and $u=(a/\bar a)^{\lfloor(s+2)/2\rfloor}$, the shift
maps $\Lambda^+$ onto $a\Lambda'^-$.
\item If an integral $h\in\U(X\oplus\langle e_0\rangle)$ has last row $\equiv(0,\dots,0,1)$ modulo $\pi$, its image in $\U(V)$
stabilises $\Lambda$; if moreover $h\equiv1$ modulo a place $w$, the image is $\equiv1$ modulo $w$.
\end{enumerate}
\end{proposition}
\lean{New.Seesaw.isUnimodular_gradedLatticeBasis_gram, New.Seesaw.gram_gradedLatticeBasis_eq_gluedGram, New.Seesaw.exists_conjTranspose_mul_gradedGram_mul_eq_signMatrix, New.Seesaw.gram_stringFrame_eq_signMatrix, New.Seesaw.conjTranspose_stringFrame_one_mul_self, New.Seesaw.stringFrame, New.Seesaw.stringSign, New.Grading.SpansLattice.isGraded, New.Grading.SpansLattice.mem_lattice_iff, New.Grading.SpansLattice.mem_latticeMinus_iff, New.Grading.SpansLattice.mem_latticePlus_iff, New.Grading.SpansLattice.stringSum_mem_latticeMinus_iff, New.Grading.SpansLattice.stringSum_mem_latticePlus_iff, New.Grading.lineIdele_inv_mul_stringSum_mem_iff, New.Grading.SpansLattice.exists_gradedLatticeBasis_mulVec_eq, New.gradedLatticeBasis_mulVec_apply_inl, New.gradedLatticeBasis_mulVec_apply_inr, New.Grading.SpansLattice.map_stringMapFin_latticePlus, New.Grading.SpansLattice.odd_of_notMem_latticeMinus, New.Grading.SpansLattice.exists_pole_of_notMem_latticeMinus, New.Grading.SpansLattice.odd_and_exists_place_of_notMem_latticeMinus, New.Grading.SpansLattice.exists_place_generates_of_notMem_latticeMinus, New.Seesaw.IsSplitSquarefree.v_algebraMap, New.Seesaw.IsSplitSquarefree.algebraMap_mem_Ô, New.Seesaw.IsSplitSquarefree.star_notMem, New.Seesaw.IsSplitSquarefree.v_apply, New.Seesaw.IsSplitSquarefree.algebraMap_eq_mul_star, New.Grading.inclXFin_mem_stdLevel, New.Grading.exists_inclXFin_sub_one_eq, New.Grading.inclXFin_mem_levelOf, New.mem_levelOf_bot}

\begin{proof}
(a) In the graded basis the block of $X$ is the identity. A string vector $b_i$ pairs only with $b_{s-1-i}$; for odd $s$ both are multiplied by $u$, and $\langle ub_i,ub_{s-1-i}\rangle=u\bar u=1$, while for even $s$ one of them is multiplied by $u/a$ and the other by $u\bar a$, and the factors $a^{-1}$ and $\bar a$ of partners cancel, so the pairing is again $1$. For odd $s$ put $v=e_0/\bar\pi+uz/N$. Then
\[
  \langle v,v\rangle=\frac1{\pi\bar\pi}+u\bar u\,\frac{-N}{N^2}=0,\qquad\langle v,uz\rangle=u\bar u\,\frac{-N}{N}=-1,\qquad\langle uz,uz\rangle=-N ,
\]
and $v$ is orthogonal to $X$ and to the other string vectors. Hence the Gram matrix of the graded basis is that of the glued form: the identity on $X$, the antidiagonal $1$ on the string with $-N$ in the middle, and for odd $s$ the block $\begin{psmallmatrix}0&-1\\-1&-N\end{psmallmatrix}$ on $(v,uz)$. Its inverse is the
glued form with the parameters $N$ and $0$ exchanged, whose entries lie in every subring containing $N$; hence the Gram matrix is hermitian and invertible over every subring. Over $\C$, the complex frame whose columns are $(\sigma_bb_b+b_{s-1-b})/\sqrt2$, with $\sigma_b=-1$ for the first $\lfloor(s+1)/2\rfloor$ indices and $+1$ for the others, and $b_m/\sqrt N$ for
the middle index $m$, has the Gram matrix $\mathrm{diag}(\sigma)$: away from the middle, the columns $b$ and $s-1-b$ span a hyperbolic plane, are orthogonal because their signs are opposite, and have norms $\pm1$, and the middle column has norm $-N/N=-1$. With the identity on $X$ this gives the signature.

(b) Every vector of the graded basis is homogeneous: the vectors of $X$ and $e_0$ have grade $0$, the string vectors their own grades, and for odd $s$ the vector $v$ combines $e_0$ with the middle vector $z$, both of grade $0$. Hence the part in each grade of a lattice vector is again a lattice vector, and $\Lambda$ is graded. A vector of $\Lambda$ lies in $X^\perp$ exactly
when its coordinates on $X\oplus\langle e_0\rangle$ vanish; as $v$ is the only basis vector with a component on $e_0$, its coefficient must vanish, and the vector is an $\hat\OO$-combination of the scaled string vectors. Hence $\Lambda^-$ is the sum of the lines $\hat\OO\ell_ib_i$, with $\ell_i$ the scale of $b_i$ in the graded basis. The projection to
$X^\perp$ kills the vectors of $X$ and sends $v$ to $uz/N$, so $\Lambda^+$ is the same sum with the line of the middle vector replaced by $\hat\OO uz/N$, i.e.\ its generator multiplied by $N^{-1}$. Let $y\in\Lambda^+\setminus\Lambda^-$. Its coordinates agree with those of an element of $\Lambda^-$ except at the middle vector, so $s$ is odd, and its middle
coordinate divided by $u$ lies in $N^{-1}\hat\OO$ but not in $\hat\OO$. Since $N=\pi\bar\pi$ with $\pi$ a uniformiser on $P$, $\bar\pi$ a uniformiser on $\bar P$ and $P\cap\bar P=\emptyset$, $N$ has valuation $1$ at the places of $P\cup\bar P$ and $0$ elsewhere; hence the middle coordinate of $y$ has a simple pole at some place $w$ above $N$. Every
$y'\in\Lambda^+$ has at most a simple pole there, so $y'-cy$ has none at $w$ for the quotient $c\in\hat\OO_w$ of the two middle coordinates: $y$ generates $\Lambda^+$ over $\Lambda^-$ at $w$.

(c) Let $\mathfrak e^i$ be the string vector $b_i$ times the element of $E$ (a power of the constant of the grading) for which the shift sends $\mathfrak e^i$ to the corresponding vector $\mathfrak e'^i$ of $Y'$ when $i<s-1$, and kills $\mathfrak e^{s-1}$. As the constant has the valuations of $a\bar a$ and
$u=(a/\bar a)^{\lfloor(s+2)/2\rfloor}$, the scale $\ell_i$ of the graded basis is $a^{s-i}\bar a^{-1-i}$ times a power of $a\bar a$, and the powers of $a\bar a$ cancel place by place. Writing $\mathfrak l_i$ for the ideal of $a^{s-i}\bar a^{-1-i}$, one finds: $\sum_it_i\mathfrak e^i\in\Lambda^-$ if and only if $t_i\in\mathfrak l_i$ in grade at least $0$ and
$t_i\in N^{-1}\mathfrak l_i$ in negative grade; and $\sum_it_i\mathfrak e^i\in\Lambda^+$ if and only if $t_i\in\mathfrak l_i$ in positive grade and $t_i\in N^{-1}\mathfrak l_i$ in grade at most $0$. For $Y'$, with $s-1$ in place of $s$, the ideals are $\mathfrak l'_i=(a^{s-1-i}\bar a^{-1-i})$, so $a\mathfrak l'_i=\mathfrak l_i$; the grade of $\mathfrak e'^i$ in $Y'$ is
$s-2-2i$, which is at least $0$ exactly when the grade $s-1-2i$ of $\mathfrak e^i$ in $Y$ is positive. Hence, for every $i<s-1$, the condition on $t_i$ for $\Lambda^+$ is the condition for $a\Lambda'^-$, and the shift, which forgets the coordinate $t_{s-1}$, maps $\Lambda^+$ onto $a\Lambda'^-$.

(d) In rational coordinates the image of $h$ acts as $h$ on $X\oplus\langle e_0\rangle$ and as $1$ on the string, so its matrix in the frame of the graded basis is $1+R^{\mathsf H}MR$, with $R$ the integral rows of the frame at $X\oplus\langle e_0\rangle$ and $M=h-1$ for even $s$; for odd $s$, because the graded basis replaces $e_0$ by $e_0/\bar\pi+uz/N$,
$M$ is $h-1$ with its last row divided by $\pi$ and its last column divided by $\bar\pi$. The integral matrices whose last row is $\equiv(0,\dots,0,1)$ modulo $\pi$ form a group, so $h^{-1}=h^{\mathsf H}$ has the same property, and its last row read in $h$ says that the last column of $h$ is $\equiv e_0$ modulo $\bar\pi$. Hence the last row of $h-1$ is divisible by
$\pi$ and its last column by $\bar\pi$; the corner is divisible by both, and at every place one of $\pi$, $\bar\pi$ is a unit, so the corner divided by $\pi\bar\pi$ is integral. Hence $M$ is integral, the image of $h$ is integral with integral inverse, and it stabilises $\Lambda$. If moreover $h\equiv1$ modulo $w$, for the places $w$ used here, which with their conjugates lie outside $P$, then $\pi$ and $\bar\pi$ are
units at $w$, $M\equiv0$ modulo $w$, and each entry of $R^{\mathsf H}MR$ is a sum of products of an entry of $M$ with two integral elements, hence $\equiv0$ modulo $w$.
\end{proof}

\begin{proposition}[The lattices of the spaces $Y_k$]\label{prop:lattice-Y}
For $k<m$ the grading of $Y_k$ carries the constant $-(\jmath M\delta_E)^{-1}$, whose valuations are those of $a\bar a$ with
$a=\alpha/\bar\beta_M$; and the point of $Y_k$ spans the graded lattice of $a$, $\pi_N$ and
$u=(a/\bar a)^{\lceil(m-k)/2\rceil}$. The idele generating $(\bar\pi_M)$ has the valuations of $\bar\pi_M$, and
$\Nm(\bar\pi_M)=M^f$.
\end{proposition}
\lean{New.Data.spansLattice_gradingY, New.Data.v_gradingY_genericConst_eq_ideleA_mul_star, New.Data.v_alpha_mul_star_eq_ideleA_norm_mul_M, New.Data.v_chosenIdele, New.Data.absNorm_piM, New.Data.gradingY, New.Data.lastGrading, New.Data.latticeScale_two_mul}

\begin{proof}
The chosen idele generates the ideal of $\bar\pi_M$ and is $1$ away from the chosen places, so it has the valuations of $\bar\pi_M$; and $\bar\pi_M$ is an integer of $K$ of norm $M$, viewed in $E$, so the ideal it generates in $E$ has norm $M^{[E:K]}=M^f$. The self-duality of $\hat\OO_E$ under $(z,u)\mapsto\psi_f(\jmath\alpha\bar\alpha zu)$ says that the order of $\jmath\alpha\bar\alpha$ at
every place is minus that of the different, which $\delta_E$ generates; so $\jmath\alpha\bar\alpha\delta_E$ is a unit at every finite place. Since $a=\alpha/\bar\beta_M$ and $\beta_M\bar\beta_M$ has the valuations of $M$, $\alpha\bar\alpha$ has the valuations of $a\bar aM$, and therefore $\jmath M\delta_Ea\bar a$ is a unit at every finite place: the constant $-(\jmath M\delta_E)^{-1}$
has the valuations of $a\bar a$. The frame of $Y_k$ was chosen (\cref{def:spaces-Y}) to be the graded basis for $a$, $\pi_N$ and $u=(a/\bar a)^{\lceil(m-k)/2\rceil}$ up to $\GL(\hat\OO)$, which is the assertion that the point spans the graded lattice; note that $u$ gains one factor $a/\bar a$ from $k=2i+2$ to $k=2i$.
\end{proof}

\subsection{Unfolding a pair of lifts}\label{ss:unfold-pair}
For data $f$ on $\Hom(V,V')$ the \emph{constant term} $C_{\mathrm{sh}}(f)(g,h)$ is the mean, over the adelic maps
$V\to V'$ of positive degree modulo the rational ones, of the theta function of the Weil translate of $f$ by $(g,h)$,
Heisenberg translated by $\mathrm{sh}$. One unfolding step (\cref{prop:B-unfolding}) writes the period of a
$\psi_U$-equivariant function $W$ on $R_V$ against a theta descent as the integral over $R_{V'}(E^+)\backslash\U(V')(\Ad)$ of
the form against
\[
  \mathcal W'(h)=\int_{[R_V]}W(r)\,C_{\mathrm{sh}}(f)(r,h)\,dr ,
\]
and $\gamma\mapsto\mathcal W'(\gamma h)$ is again $\psi_U$-equivariant.
\lean{New.Grading.constTerm, New.Grading.constTermAtStringMap, New.Grading.unfoldStep, New.Seesaw.constTermAtStringMap_unipotent_mul_left, New.Seesaw.constTermAtStringMap_unipotent_mul_right, New.Seesaw.mul_constTermAtStringMap_unipotent_mul}

A \emph{pair of edges} $Y\to Y'\to Y''$ consists of gradings of $V$, $V'$, $V''$ with strings of lengths $s$, $s-1$
and $s-2$ (or $0$ if $s=1$), and an identification of the part of grade zero of $V^-\oplus V''$, where $V^-$ is $V$ with
the form negated, with $X^-\oplus X''$ plus $p$ hyperbolic planes. It gives a splitting of $\Hom(V',V^-\oplus V'')$ down to
$\Hom(X',X^-\oplus X'')$ and an embedding of $\U(X_0)$ into $\U(X'')$. The \emph{middle} pairs have one plane, the
\emph{last} pair none. The \emph{join} of data $f$, $f'$ on the two edges is the datum $f\star f'$ on
$\Hom(V',V^-\oplus V'')$ whose theta value is $z\ne0$ times the product of those of $f$ and $f'$.
\lean{New.Seesaw.DoubleLiftData, New.Seesaw.HasConsecutiveStrings.c_eq, New.Seesaw.HasConsecutiveStrings.length_last, New.Seesaw.HasConsecutiveStrings.odd_length, New.Seesaw.HasConsecutiveStrings.length_middle, New.Seesaw.DoubleLiftData.toHasConsecutiveStrings, New.Seesaw.GradeZeroPart, New.Seesaw.GradeZeroPart.splitting, New.Seesaw.GradeZeroPart.hom, New.Seesaw.GradeZeroPart.index, New.Seesaw.GradeZeroPart.incl, New.Seesaw.GradeZeroPart.gram_incl, New.Seesaw.GradeZeroPart.incl_eq_zero, New.Seesaw.GradeZeroPart.mem_range_index, New.Seesaw.GradeZeroPart.grade_add_eq_zero, New.Seesaw.GradeZeroPart.J_eq_zero_of_add_ne_zero, New.Seesaw.GradeZeroPart.proj_pos_mul_J, New.Seesaw.GradeZeroPart.inv_J_apply_eq_zero_of_pos, New.Seesaw.GradeZeroPart.card_pos_eq_card_neg, New.Seesaw.DoubleLiftData.homZeroPart, New.Seesaw.DoubleLiftData.splitting, New.Seesaw.DoubleLiftData.zeroPart, New.Seesaw.DoubleLiftData.embX, New.Seesaw.DoubleLiftData.det_embX, New.Seesaw.DoubleLiftData.intertwines, New.Seesaw.DoubleLiftData.gradeOne, New.Seesaw.DoubleLiftData.even_length_middle, New.Seesaw.DoubleLiftData.even_gradeOf, New.Seesaw.DoubleLiftData.even_card_add, New.Seesaw.middleDoubleLift, New.Seesaw.lastDoubleLift, New.Seesaw.stackedStringMap, New.Seesaw.stackedStringMapFin, New.Seesaw.DoubleLiftData.stackedStringMap_apply}

\begin{theorem}[Two unfolding steps are one double lift]\label{thm:two-steps}
Let $Y\to Y'\to Y''$ be a pair of edges with definite parts $X$, $X'$, $X''\supseteq X_0$. There is $z\in\C$ such that for
all characters $a$, $b$ of $\U(1)(\Ad)$, data $f$, $f'$ on the two edges, points $k\in\U(X_0)(\Ad)$ and continuous,
rationally invariant, $\psi_U$-equivariant $W$ on $R_V$,
\begin{multline*}
  \int_{[R_V]}\int_{[R_{V'}]}W(g)\,C_{\mathrm{sh}}(f)(g,r)\,C_{\mathrm{sh}'}(f')(r,k)\,dr\,dg\\
  =\int_{[\U(X)]}W(g)\,z\int_{[\U(X')]}\theta\Bigl(\omega_{a,b}(g,r,k)\,Z\bigl(\varrho(\mathrm{sh}^*,\mathrm{sh}')(f\star f')\bigr)\Bigr)\,dr\,dg,
\end{multline*}
where $\varrho(\mathrm{sh}^*,\mathrm{sh}')$ is the Heisenberg translation by the stacked map $V'\to V^-\oplus V''$, $Z$ is the zero
term along the splitting, and $\omega_{a,b}$ is the joint Weil action of $\U(X)\times\U(X')\times\U(X'')$. The zero term
intertwines the joint Weil actions of the graded and the definite spaces with no further scalar.
\end{theorem}
\lean{New.Seesaw.DoubleLiftData.exists_integral_constTermAtStringMap_mul_eq, New.Seesaw.DoubleLiftData.constTerm_mul_constTerm_eq_inv_mul_zeroTerm, New.Seesaw.DoubleLiftData.zeroTerm_stackedStringMap_planes, New.Seesaw.DoubleLiftData.exists_add_eq_homAtPoint_stackedStringMap, New.Seesaw.formAdjoint_stackedStringMap_mul_stackedStringMap, New.Seesaw.DoubleLiftData.zeroTerm_splitting_tripleWeil, New.Seesaw.GradeZeroPart.zeroTerm_hom_omega, New.Seesaw.integrable_algFormCoeff_mul_theta_tripleWeil, New.Seesaw.arrowProduct_integrated_polyVector}

\begin{proof}
\emph{The product of two constant terms.} The constant term of $f$ at $w$ is the mean of $\theta(\varrho(w+a)f)$ over the maps $a\colon V\to V'$ of positive degree, adelic modulo rational, and similarly for $f'$ at $w'$ with $a'\colon V'\to V''$. A map $V'\to V^-\oplus V''$ of positive degree for the grading of $V^-\oplus V''$ is a pair
$(a^*,a')$ with $a$, $a'$ of positive degree, so the mean over such maps is the iterated mean over $a$ and $a'$; and the theta value of the join is $z$ times the product of the theta values. Hence the product of the two constant terms is $z^{-1}$ times the value at $0$ of the zero term, along the splitting that removes the maps of nonzero degree, of
the join of the two Heisenberg translates.

\emph{The stacked map.} Write $T=\mathrm{sh}\colon V\to V'$ and $T'=\mathrm{sh}'\colon V'\to V''$ and $S=(T^*,T')$. Its adjoint for $V^-\oplus V''$ is $(-T,T'^*)$, so $S^\dagger S=-TT^*+T'^*T'$; both products depend only on the grading of $V'$ and the common constant, and they agree, so $S^\dagger S=0$: $S$ is isotropic, and translations by $S$ and by
any multiple of $S$ compose without a phase.

\emph{The unipotent variables.} Let $n$ be a unipotent point of $\U(V')$. In rational coordinates $S(n-1)$ is the sum of a map of positive degree and a map of degree zero that sends $X'$ into an isotropic half of the hyperbolic planes of grade zero; the matrix of the latter is the constant of the grading times the block of $n-1$ with rows the basis
vectors of grade one and columns $X'$. The zero term ignores translations by maps of positive degree, and translating by the stacked map and then by $S(n-1)$ is translating by $Sn$; so $z^{-1}$ times the zero term of the translate by $S$ of the join, evaluated at the map of degree zero just described, is the product of the constant terms of $f$ at
$n^{-1}T$ and of $f'$ at $T'n$. A unipotent element of $R_V$ in the first argument of the constant term moves into the map on the right, and one of $R_{V'}$ in the second argument moves into it on the left; this is what the last sentence says.

\emph{The integrals.} Write $g=u\,\iota_X(g_0)$ and $r=u'\iota_{X'}(r_0)$ with unipotent $u$, $u'$ and integrate first over the unipotent groups, whose quotients by the rational points have volume $1$. Over $[U_V]$ the integrand is invariant: $W$ transforms by $\psi_U$, and the product of the constant term with the characters of the pair transforms by
$\psi_U^{-1}$, because the string of $V$ has odd length and the constants of the two gradings agree. Over $[U_{V'}]$, by the previous paragraph, the integrand is the zero term evaluated at the degree-zero map attached to $u'$; as $u'$ runs over $[U_{V'}]$, these maps run over the adelic maps into the isotropic half of the planes modulo the rational ones,
with a fixed constant Jacobian, so the mean over $[U_{V'}]$ turns the zero term along the maps of nonzero degree into the zero term $Z$ along the splitting that removes also the hyperbolic planes. Collecting the constants of the two unfoldings and $z^{-1}$ into one constant $z$ gives the formula, with the remaining action of $(g_0,r_0,k)$ on
the translated datum.

\emph{No further scalar.} The zero term is taken in two stages: first to the part of grade zero of $\Hom(V',V^-\oplus V'')$, which is $\Hom(X',Z_0)$ with $Z_0$ the grade-zero part of $V^-\oplus V''$, and then off the hyperbolic planes of $Z_0$. At each stage the law of the zero term along a splitting
(\cref{ss:B-siegel-weil}) carries the Weil operator of the Levi component to that of the smaller space, times a power of the splitting character $\upsilon$ at the determinant. The characters of the pair are the determinant characters $a$, $b$, twisted by powers of $\upsilon$ depending on the ranks; on the graded spaces the twists of the target characters are
$\upsilon^{\pm1}$ according to the parities of the ranks, on the definite spaces they are trivial, and in the product of the two target characters $\upsilon$ cancels against the powers produced by the two stages, because the length of the string of $V'$ and the sum of the ranks of $V$ and $V''$ are even. Hence the zero term carries the joint Weil
operator of $(\iota g_0,\iota r_0,\iota k)$ to that of $(g_0,r_0,k)$ with no scalar.
\end{proof}

The \emph{finite evaluation} of a pair of edges is the linear map on joined finite data which translates by the finite
part of the stacked map, takes the finite part $Z_f$ of the zero term along the splitting, and splits the result into a
tensor of finite data for $\Hom(X,X')$ and $\Hom(X',X'')$ (a left inverse of the join). For the pair
$Y_{2i}\to Y_{2i+1}\to Y_{2i+2}$ it is the evaluation of the middle pair for $i<\varrho$ and of the last pair for
$i=\varrho$. The finite part $Z_f$ is not zero at the vacuum, obeys the law of the zero term for finite points of the
parabolic subgroup, and is equivariant for the finite joint Weil action.
\lean{New.Seesaw.DoubleLiftData.finEval, New.Seesaw.DoubleLiftData.arrowTensor_finEval, New.Seesaw.DoubleLiftData.finEval_arrowProduct_equivariant, New.Seesaw.zeroTermFin, New.Seesaw.zeroTermFin_eq, New.Seesaw.exists_zeroTermFin_ne_zero, New.Seesaw.zeroTermFin_omega, New.Seesaw.zeroTermFin_mem_finitePart, New.Seesaw.zeroTermFin_vacElem_ne_zero, New.Seesaw.exists_zeroTerm_rho_integrated_eq, New.Space.Splitting.iFin, New.Space.Splitting.isotropicBasisFin, New.Space.Splitting.conjTranspose_isotropicBasisFin, New.Data.doubleLiftEvalY, New.Data.doubleLiftEvalY_inclXFin}

These properties of $Z_f$ are proved as follows. The finite part is $f\mapsto Z(v\otimes f)$ for the Fock vacuum $v$ moved by the real part of a point adapted to the splitting, and the law of $Z$ for the parabolic subgroup restricts to finite points, which do not move $v$. At the image $\omega(1,g_f)\mathrm{vac}$ the zero
term is that of $\omega(g)\mathrm{vac}$, a nonzero multiple of the vacuum; so $Z_f$ does not vanish on the finite part, and since the finite translates of the vacuum span the finite part and $Z_f$ commutes with them up to a root of unity, it does not vanish at the vacuum. For the equivariance, the images
$\iota_X$, $\iota_{X'}$, $\iota_{X''}$ fix the shifts (they act on $X$, $X'$, $X''$, which the shifts kill or avoid), so the joint action passes the translation by the stacked map, and by the last part of the proof of \cref{thm:two-steps} it passes the zero term with no scalar; it then splits over the two factors of the join.

\subsection{The unfolded functions}\label{ss:unfolded}
Fix a source form $\varphi$, a linear form $t$ on its values and an archimedean point $g$ of $\U(Y_0)$. The
\emph{outer variables} are points $g_j=(\xi_j,y_j)$ of $\U(Y_j)(\Ad)$, $\xi_j$ archimedean and $y_j$ finite. The datum of
lift $j+1$ moved by $g_j$ is $\omega(g_j,1)(\varphi_{\rm harm}\otimes\Upsilon_{j+1})$, and the \emph{unfolded functions} are
defined by
$\mathcal W_0=\mathcal W_\varphi$ on $R_0$ and
\[
  \mathcal W_{k+1}(h)=\int_{[R_k]}\mathcal W_k(r g_k)\,C_{\mathrm{sh}}\bigl(\omega(g_k,1)(\varphi_{\rm harm}\otimes\Upsilon_{k+1})\bigr)(r,h)\,dr
\]
for $k\le2\varrho+1$, where $\mathcal W_k(rg_k)$ is read as $\mathcal W_0(r)$ for $k=0$. Each $\mathcal W_{k+1}$ is continuous, left
invariant under $R_{k+1}(E^+)$, and $\gamma\mapsto\mathcal W_{k+1}(\gamma g_{k+1})$ is $\psi_U$-equivariant on $R_{k+1}$. The
\emph{translated data} of lift $k+1$ are the finite data $\omega_f(y_k,y_{k+1})\Upsilon_{k+1}$, and the \emph{evaluated datum}
$\phi_i(y)$ of the pair $i$ is the finite evaluation of the two translated data of lifts $2i+1$, $2i+2$.
\lean{New.Data.movedDatum, New.Data.twistedPairWeil_movedDatum, New.Data.constTermAtStringMap_movedDatum, New.Data.unfoldStageOn, New.Data.unfoldStage, New.Data.unfoldStageOn_succ, New.Data.unfoldStageOn_update, New.Data.unfoldStage_update, New.Data.isEquivariant_unfoldStageOn, New.Data.continuous_unfoldStage, New.Data.fixesLeft_ratBessel_unfoldStage, New.Data.integral_unfoldStageOn, New.Data.translatedPointDatumY, New.Data.evaluatedDatum, New.Data.evaluatedDatum_out, New.Data.liftWeilFin_pointDatumY, New.Data.liftWeilFin_translatedPointDatumY, New.Data.map_evaluatedDatum_of_conj_mem, New.Data.inclXFin_mem_chosenLevel}

Each of these properties is the conclusion of \cref{prop:B-unfolding} at the step $k+1$, applied to the function $\gamma\mapsto\mathcal W_k(\gamma g_k)$, which is $\psi_U$-equivariant by the previous step; the constants of the gradings of all $Y_k$ are the same number $-(\jmath M\delta_E)^{-1}$, as the proposition requires.
By construction $\mathcal W_{k+1}$ depends only on the outer variables $g_j$ with $j\le k$.

\begin{theorem}[The unfolded functions are upper chains]\label{thm:unfolded-upper}
Let $\Theta_i(y)$ be the composite of the first $i$ double lifts $X_0\to\dots\to X_{2i}$ with kernels of \cref{ss:composite-lift}
and finite data $\phi_0(y),\dots,\phi_{i-1}(y)$. For every $i$ with $2i\le2\varrho+1$ and all archimedean $\xi$ there is a linear
form $l$, independent of $\vartheta$, $\varphi$ and $y$, such that $\mathcal W_{2i}$ restricted to $R_{2i}$ is the Bessel function
$\mathcal W$ of $\Theta_i(y)\varphi$ at $l$; and for $2i<2\varrho+1$ and every archimedean $x$ there is $l$ with
\[
  \mathcal W_{2i+2}\bigl(\iota_X(g)\,(x,y_{2i+2})\bigr)=l\bigl((\Theta_{i+1}(y)\varphi)(\iota g)\bigr)\qquad(g\in\U(X_{2i+2})(\Ad)).
\]
\end{theorem}
\lean{New.Data.exists_unfoldStageOn_even_eq_sourceFn, New.Data.exists_unfoldStage_eq_algFormCoeff, New.Data.exists_unfoldStage_eq_algFormCoeff_of_sourceFn, New.Data.exists_zeroTerm_eq_integrated, New.Seesaw.integral_theta_tripleWeil_eq_lift_comp}

\begin{proof}
Induction on $i$. For $i=0$ the function on $R_0$ is $\mathcal W_\varphi$, the Bessel function of $\varphi=\Theta_0(y)\varphi$ at $t$. Suppose that the function that $2i$ steps leave on $R_{2i}$ is the Bessel function of $\psi=\Theta_i(y)\varphi$ at a linear form $l$, that is, the $\psi_U$-equivariant function equal to the
coefficient $l(\psi(g))$ on $\iota\U(X_{2i})(\Ad)$.

\emph{Two steps.} Steps $2i+1$ and $2i+2$ are the left side of \cref{thm:two-steps} for the pair $Y_{2i}\to Y_{2i+1}\to Y_{2i+2}$, with $W$ this function, $f$ and $f'$ the data of lifts $2i+1$ and $2i+2$ moved by the outer variables, and the third point $(x,y_{2i+2})$. By the theorem they equal $z$ times a
double theta integral over $[\U(X_{2i})]\times[\U(X_{2i+1})]$ of the datum $Z(\varrho(\mathrm{sh}^*,\mathrm{sh}')(f\star f'))$ moved by the joint Weil action.

\emph{The datum is a pure tensor.} Each of the operations (moving by the outer variables, the join, the translation by the stacked map and the zero term) acts separately on the real and on the finite factor of a pure tensor, and on the real factor the characters act through their weights, which do not depend on $\vartheta$. Hence the datum is
$v\otimes\phi_i(y)$, carried to $\Hom(X_{2i+1},X_{2i}^-\oplus X_{2i+2})$ by the join, with a Fock vector $v$ that depends on the archimedean outer variables and on $x$ but not on $\vartheta$, $\varphi$ or $y$; the finite factor is the finite evaluation of the two translated data, by definition the evaluated datum $\phi_i(y)$.

\emph{Double lifts.} It remains to see that, for every Fock vector $u$ and linear form $l$, there is $l''$ such that the double integral of $\theta$ of $u\otimes\Phi$ moved by $(g_0,r,k)$ against the coefficient $l(\psi(g_0))$ is the coefficient at $l''$ of the composite of the two lifts with finite data $\Phi$ applied to $\psi$. For $u=PQ$ a
product of a polynomial on $\Hom(X_{2i},X_{2i+1})$ and one on $\Hom(X_{2i+1},X_{2i+2})$ (the variables of $P$ transposed into the block of $X_{2i}$), the datum $PQ\otimes\Phi$ is the join of $P\otimes\Phi_1$ and $Q\otimes\Phi_2$ for $\Phi=\Phi_1\otimes\Phi_2$, and its kernel is $z$ times the product of the two theta kernels; so the double integral is
an iterated theta lift between definite spaces, which \cref{sec:kernels} computes as a class sum of jets, and $l''$ exists. Sums of such products are the Fock polynomials; these are dense in the Fock space, the integral is continuous in $u$, and the families given by linear forms $l''$ make up a finite-dimensional, hence closed, space; so the
claim holds for every $u$, and with $u=v$ it gives the second display. As $\mathcal W_{2i+2}$ is $\psi_U$-equivariant on $R_{2i+2}$, it is determined by its values on $\iota\U(X_{2i+2})(\Ad)$, and the first statement for $i+1$ follows with $l''$ in place of $l$.
\end{proof}

\subsection{The support of the unfolded functions}\label{ss:support}
A finite point $x$ of $\U(Y_j)$ \emph{has the lattice bounds} for $(\mathfrak q,k)$ if $\bar{\mathfrak q}^k\Lambda^-\subseteq\Lambda^-_x$ and
$\mathfrak q^k\Lambda^-_x\subseteq\Lambda^-$. It has \emph{bounded orbit} if some point $rx$, $r\in R_j(\Ad_{\mathrm f})$, has a graded
lattice and the bounds. The lattice of $x$ \emph{splits} along $X$ at a place $w$ if there $\Lambda^+_x\subseteq\Lambda^-_x$.
\lean{New.Grading.HasLatticeBounds, New.Grading.HasBoundedOrbit, New.Grading.LatticeSplitsAt, New.Grading.LatticeSplitsAt.conj, New.Grading.LatticeSplitsAt.idem_smul_sub_stringProj_mem}

\begin{lemma}[The shift at a nonzero constant term]\label{lem:frequency}
Let the constant term at $\mathrm{sh}$ of the Weil translate by $(x,x')$ of a datum fixed by the Heisenberg translations
by $b\,\hat\OO^{\iota'\times\iota}$ be nonzero. Then some matrix $T$ with $T\Lambda_{x}\subseteq\Lambda_{x'}$ agrees with $(\bar b/\alpha)\,\mathrm{sh}$ up to
a matrix that does not lower grades; and if $r$, $r'$ are finite points of the Bessel groups for which $\Lambda_{rx}$, $\Lambda_{r'x'}$
are graded, then $(\bar b/\alpha)\,\mathrm{sh}\,\Lambda^+_{rx}\subseteq\Lambda'^-_{r'x'}$.
\end{lemma}
\lean{New.Grading.exists_mapsTo_lattice_hasDegreeGe_sub_smul_stringMap, New.Grading.smul_map_stringMapFin_latticePlus_le, New.Grading.IsGraded.mulVec_mem_of_homogeneous, New.Grading.zeroTermFin_finHeisenbergOp_stringMap_ne_zero, New.Grading.xFin_mulVec_finCoord_homAtPoint_sub, New.Grading.hasDegreeGe_zero_mul_mul_sub, New.Grading.hasDegreeGe_zero_ratCoordMatrixFin}

\begin{proof}
\emph{The finite zero term.} The constant term at $\mathrm{sh}$ of the translate of $t\otimes f$ by $(x,x')$ is the theta value of the zero term, along the splitting $D$ that removes the maps of positive degree, of the translate by $\mathrm{sh}$ (\cref{ss:B-unfolding}); each operation acts separately on the real and finite factors, and the characters contribute
a scalar. Hence if the constant term does not vanish, the finite zero term $Z_f$ does not vanish at $\varrho(\mathrm{sh}_f)\omega_f(x,x')f$.

\emph{The support of $Z_f$.} Let $k$ be the finite isotropic frame of $D$ and $k^\perp$ the kernel of $k^{\mathsf H}$. Suppose $g$ is fixed by the Heisenberg operators of $\mathcal M\hat\OO^n$, with $\mathcal M\mathcal M'^{\mathsf H}=1$, and $Z_f(\varrho(y)g)\ne0$. Rescale $k$ by a $t$ so that $\mathcal M'^{\mathsf H}kt$ is integral with an integral left inverse $L$. A
vector $\zeta$ of the isotropic subspace is a period of $Z_f$, and if it is also a period of $g$, then the commutator of the translations by $\zeta$ and by $y$, the scalar $\psi(\langle\zeta,y\rangle)$, fixes the nonzero vector $Z_f(\varrho(y)g)$, hence equals $1$. For $\xi$ integral, $\zeta=kt\,\xi$ is such a vector: it is
$\mathcal M(\mathcal M'^{\mathsf H}kt\,\xi)$ with $\mathcal M'^{\mathsf H}kt\,\xi$ integral. Hence $\psi(\langle\xi,(kt)^{\mathsf H}y\rangle)=1$ for all integral $\xi$, and $v=(kt)^{\mathsf H}y$ is integral, since the standard lattice is its own dual. With $l=L^{\mathsf H}v$ one has $(kt)^{\mathsf H}\mathcal M'l=v$, so $y-\mathcal M'l\in k^\perp$.

\emph{The matrix $T$.} Apply this to $g=\omega_f(x,x')f$, which is fixed by the translations by $b\,\omega_f(x,x')\hat\OO^{\iota'\times\iota}$, so that $\mathcal M'$ is $\bar b^{-1}$ times the action of $(x,x')$; and $k^\perp$ consists of the maps of degree at least $0$. Hence the finite coordinates of $\mathrm{sh}$, that is $\mathrm{sh}_f/\alpha$, are
$\bar b^{-1}$ times the image under $(x,x')$ of an integral matrix $l$, up to a map of degree at least $0$. In rational coordinates the image of $l$ is $T=x'lx^{-1}$, which maps $\Lambda_x$ into $\Lambda_{x'}$, and $T-(\bar b/\alpha)\mathrm{sh}$ does not lower grades.

\emph{Graded points.} For $r$, $r'$ in the finite Bessel groups, $r'Tr$ maps $\Lambda_{rx}$ into $\Lambda_{r'x'}$; the matrices of $r$ and $r'$ in rational coordinates do not lower grades (they fix the string vectors of grade $\ge0$ up to vectors of higher grade), degrees add under multiplication, and $\mathrm{sh}$ is homogeneous of degree $-1$, so
$r'Tr-(\bar b/\alpha)\mathrm{sh}$ again does not lower grades. Its part of degree $-1$ is therefore $(\bar b/\alpha)\mathrm{sh}$. If $\Lambda_{rx}$ and $\Lambda_{r'x'}$ are graded, the part of degree $-1$ of a map between them maps the first into the second, being the sum over the grades $d$ of the parts from grade $d$ to grade $d-1$, each of which
maps the graded pieces correctly. Finally $\mathrm{sh}$ kills $X$, so $\mathrm{sh}\,y=\mathrm{sh}\,y^+$ for the projection $y^+$ of $y$ to $X^\perp$, and $\mathrm{sh}$ has image in $X'^\perp$. Hence $(\bar b/\alpha)\mathrm{sh}\,\Lambda^+_{rx}=(\bar b/\alpha)\mathrm{sh}\,\Lambda_{rx}\subseteq\Lambda_{r'x'}\cap X'^\perp=\Lambda'^-_{r'x'}$.
\end{proof}

\begin{proposition}[The sandwich]\label{prop:sandwich}
Let $Y$, $Y'$ be gradings with strings of lengths $s$, $s-1$ whose points span graded lattices, with $\mathrm{sh}\Lambda^+=a\Lambda'^-$.
\begin{enumerate}[label=\textup{(\alph*)}]
\item If $\mathfrak q^k\Lambda^-_x\subseteq\Lambda^-$ and $(\bar{\mathfrak q}/a)\,\mathrm{sh}\,\Lambda^+_x\subseteq\Lambda'^-_{x'}$, then $\bar{\mathfrak q}^{k+1}\Lambda'^-\subseteq\Lambda'^-_{x'}$.
\item Let $\mathfrak q$ be a unit above $N$. If $\bar{\mathfrak q}^k\Lambda^-\subseteq\Lambda^-_x$, then either $x$ has the lattice bounds for
$(\mathfrak q,k)$, or $s$ is odd and the lattice of $x$ splits at some place dividing $\pi_N$.
\item If the lattice of $x$ splits at $w$, the integral scalars of $X$ supported at $w$ and $\bar w$ preserve $\Lambda_x$.
\item If every point $rx$ with graded lattice has the bounds, $x$ has bounded orbit; if some $rx$ is integral, $x$ has
bounded orbit for $k=0$.
\end{enumerate}
\end{proposition}
\lean{New.Grading.smul_latticeMinus_le_of_map_latticePlus, New.Grading.SpansLattice.hasLatticeBounds_or_latticeSplitsAt, New.Grading.latticeSplitsAt_of_generates, New.Grading.LatticeSplitsAt.conj_scalarFin_mem, New.Grading.inclXFin_scalarFin_mulVec, New.Grading.inclXFin_scalarFin_map_evalAt, New.Grading.HasBoundedOrbit.of_forall, New.Grading.SpansLattice.hasBoundedOrbit_zero, New.Grading.exists_isGraded}

\begin{proof}
\emph{Duality.} The lattice $\Lambda_x$ is self-dual and $X^\perp$ is nondegenerate, so the dual of $\Lambda^-_x=\Lambda_x\cap X^\perp$ in $X^\perp$ is the projection $\Lambda^+_x$, and conversely. The dual of $cL$ is $\bar c^{-1}L^\vee$, and duality reverses inclusions. Hence an inclusion $\mathfrak q^k\Lambda^-_x\subseteq\Lambda^-$ is equivalent to
$\bar{\mathfrak q}^k\Lambda^+\subseteq\Lambda^+_x$, and $\bar{\mathfrak q}^k\Lambda^-\subseteq\Lambda^-_x$ is equivalent to $\mathfrak q^k\Lambda^+_x\subseteq\Lambda^+$.

(a) By duality $\bar{\mathfrak q}^k\Lambda^+\subseteq\Lambda^+_x$. Apply $(\bar{\mathfrak q}/a)\mathrm{sh}$ and use $\mathrm{sh}\Lambda^+=a\Lambda'^-$:
\[
  \bar{\mathfrak q}^{k+1}\Lambda'^-=(\bar{\mathfrak q}/a)\,\mathrm{sh}\,\bar{\mathfrak q}^k\Lambda^+\subseteq(\bar{\mathfrak q}/a)\,\mathrm{sh}\,\Lambda^+_x\subseteq\Lambda'^-_{x'} .
\]

(b) By duality $\mathfrak q^k\Lambda^+_x\subseteq\Lambda^+$, hence $\mathfrak q^k\Lambda^-_x\subseteq\mathfrak q^k\Lambda^+_x\subseteq\Lambda^+$. If $\mathfrak q^k\Lambda^-_x\subseteq\Lambda^-$, the bounds hold. Otherwise take $y\in\Lambda^-_x$ with $\mathfrak q^ky\in\Lambda^+\setminus\Lambda^-$. By \cref{prop:graded-lattice}(b), $s$ is odd and $\mathfrak q^ky$ generates
$\Lambda^+$ over $\Lambda^-$ at a place $w$ above $N$. There $\mathfrak q$ and $\bar{\mathfrak q}$ are units, so the two inclusions read $\Lambda^-\subseteq\Lambda^-_x$ and $\Lambda^+_x\subseteq\Lambda^+$ at $w$, and
\[
  \Lambda^+_x\subseteq\Lambda^+=\Lambda^-+\hat\OO_w\,\mathfrak q^ky\subseteq\Lambda^-_x\qquad\text{at }w ,
\]
because $y\in\Lambda^-_x$: the lattice of $x$ splits at $w$. Since $\Lambda_x$ is self-dual, splitting at $w$ is equivalent to splitting at $\bar w$, so the place may be taken in $P$; and it divides $\pi_N$.

(c) The point of $\U(X^\perp\oplus X)$ that scales $X$ by $z$ sends $y$ to $y+(z-1)(y-y^+)$, with $y^+$ the projection to $X^\perp$. If the lattice splits at $w$, then for $y\in\Lambda_x$ the component at $w$ of $y^+$ lies in $\Lambda^-_x\subseteq\Lambda_x$, hence so does that of $y-y^+$; the same holds at $\bar w$, as splitting at $w$ implies splitting at $\bar w$. For $z$
integral, of norm one and equal to $1$ off $w$ and $\bar w$, the vector $(z-1)(y-y^+)$ is supported at $w$, $\bar w$ and is there an integral multiple of a lattice vector; so the scaling maps $\Lambda_x$ into itself, and so does its inverse, the scaling by $\bar z$. In other words, its conjugate by $x$ is integral and trivial off $w$ and
$\bar w$.

(d) Every orbit of the finite Bessel group contains a point whose lattice is graded, by the Iwasawa decomposition for the flag of the grades (\cref{ss:B-unfolding}); if all such points have the bounds, one of them has both properties, which is the definition of a bounded orbit. If $rx$ is integral, its lattice is $\Lambda$, which is graded by
\cref{prop:graded-lattice}(b) and has the bounds for $k=0$.
\end{proof}

\begin{theorem}[Support of the unfolded functions]\label{thm:support}
Let $\varphi$ be ramified and $y_0\in J_M$.
\begin{enumerate}[label=\textup{(\alph*)}]
\item If the lattice of $x$ splits along $X_{2j+2}$ at a place $w\mid\pi_N$, or if $x^{-1}z_wx\in J_{2j+2}$ for the scaling $z_w$
of \cref{ss:composite-lift} acting on $X_{2j+2}$, then $\mathcal W_{2j+2}(\gamma(a,x))=0$ for all $\gamma\in R_{2j+2}(\Ad)$ and
archimedean $a$.
\item If $\mathcal W_j(a,x)\ne0$ for some $j=k+1\le2\varrho+1$ and $r'\in R_j(\Ad_{\mathrm f})$ makes $\Lambda_{r'x}$ graded, then $r'x$ has the
lattice bounds for $(\mathfrak q,j)$. In particular $x$ has bounded orbit.
\end{enumerate}
\end{theorem}
\lean{New.Data.unfoldStage_eq_zero_of_latticeSplitsAt, New.Data.unfoldStage_eq_zero_of_scalingGenerator, New.Data.hasLatticeBounds_of_unfoldStage_ne_zero, New.Data.finHeisenbergOp_generator_smul_pointDatumY, New.Data.map_stringMapFin_latticePlus_gradingY}

\begin{proof}
(a) \emph{The scaling case.} Put $x$ in the last outer variable: by \cref{thm:unfolded-upper}, with $y_{2j+2}=x$, the value $\mathcal W_{2j+2}(\gamma(a,x))$ is a coefficient of $(\Theta_{j+1}(y)\varphi)(g)$ for a point $g$. The image $Z$ of the scalar $z_w$ on $X_{2j+2}$ lies in the Bessel group, and
$x^{-1}Zx\in J_{2j+2}$ means that $(1,Z)$ fixes the translated datum of lift $2j+2$, since the level fixes the datum $\Upsilon_{2j+2}$ and the tame characters are trivial on it. The finite evaluation carries $\omega(1,Z)$ on the data of the $Y_k$ to $\omega(1,z_w)$ on the second factor of $\phi_j(y)$
(\cref{ss:unfold-pair}), so $\omega(1,z_w)$ on the last factor fixes the tensor of the chain $\Theta_{j+1}(y)$. By the argument of \cref{prop:characters-N}, the chain then sends $\psi-z_w\psi$ to $0$ for every $\psi$, and since $\varphi$ is ramified, $\varphi$ is a sum of such differences; so the value vanishes.

\emph{The splitting case.} If the lattice of $x$ splits at $w\mid\pi_N$, \cref{prop:sandwich}(c) shows that $x^{-1}z_wx$ is integral and trivial off $w$ and $\bar w$. The places $w$, $\bar w$ lie above $N$, so they are neither chosen places nor places of $Q$ (all places above $N$ are split and lie in $P\cup\bar P$, which $Q$ avoids); so
$x^{-1}z_wx$ reduces to $1$ at all level places and lies in $J_{2j+2}$, and the first case applies.

(b) Induction on $k$, for all points and all graded translates. Suppose $\mathcal W_{k+1}(a,x)\ne0$. By the definition of $\mathcal W_{k+1}$ as an integral over $[R_k]$, some $r\in R_k$ gives a nonzero integrand: the function of the previous stage is nonzero at $rg_k$ and the constant term of the moved datum of lift $k+1$ is nonzero at
$(rg_k,(a,x))$. Let $r_0$ in the finite Bessel group make the lattice of $r_0ry_k$ graded. For $k=0$ take $r_0=r^{-1}$: the point $y_0\in J_M$ is integral, so its lattice is $\Lambda$, which is graded and has the bounds for $0$ (\cref{prop:sandwich}(d)); for $k>0$ the function of the previous stage at $rg_k$ is a value of $\mathcal W_k$,
so the induction hypothesis gives the bounds for $(\mathfrak q,k)$ at $r_0ry_k$.

\emph{Frequency.} Let $b=\mathfrak q c$, with $c$ the chosen idele, generate the modulus of the twist, the product of the chosen places and the places of $Q$. The datum of lift $k+1$ is fixed by the Heisenberg translations by $b\hat\OO$: $\Upsilon_{k+1}$ is the vacuum away from the level places and the indicator of the inclusion matrix modulo the
modulus there, a combination of translates $\varrho(t)\mathrm{vac}$ with $\bar bt$ integral, and the translations by $b\hat\OO$ commute with these and fix the vacuum. With $a=\alpha/\bar c$ one has $\bar b/\alpha=\bar{\mathfrak q}/a$. Hence \cref{lem:frequency} gives $(\bar{\mathfrak q}/a)\,\mathrm{sh}\,\Lambda^+_{r_0ry_k}\subseteq\Lambda'^-_{r'x}$.

\emph{Sandwich.} By \cref{prop:lattice-Y} and \cref{prop:graded-lattice}(c), $\mathrm{sh}\Lambda^+=a\Lambda'^-$. With the upper bound $\mathfrak q^k\Lambda^-_{r_0ry_k}\subseteq\Lambda^-$, \cref{prop:sandwich}(a) gives the lower bound $\bar{\mathfrak q}^{k+1}\Lambda'^-\subseteq\Lambda'^-_{r'x}$. The idele $\mathfrak q$ is a unit above $N$, since no place of
$Q$ lies above $N$; so \cref{prop:sandwich}(b) gives the upper bound unless the string of $Y_{k+1}$ has odd length, i.e.\ $k+1$ is even, and the lattice of $r'x$ splits at a place dividing $\pi_N$. In that case (a) gives $\mathcal W_{k+1}(\gamma(a,r'x))=0$ for all $\gamma\in R_{k+1}(\Ad)$, and with $\gamma=(1,r'^{-1})$ this contradicts $\mathcal W_{k+1}(a,x)\ne0$.
Hence $r'x$ has the bounds for $(\mathfrak q,k+1)$. As every graded point of the orbit of $x$ has the bounds, $x$ has bounded orbit by \cref{prop:sandwich}(d).
\end{proof}

\subsection{Cuts}\label{ss:cuts}
Let $S$, $A\supseteq J$ be subgroups of a locally compact group $H$, and $T$ a finite set of cosets of $J$. The \emph{cover
multiplicity} of $t\in H/J$ is the number of classes $r\in S/(S\cap A)$ with $r^{-1}t\in T$ (zero if there are infinitely
many). A \emph{cut} for stage $j>0$ is a finite set $\Omega_j$ of classes of $\U(Y_j)(\Ad_{\mathrm f})$ modulo $J_j$ which is stable
under $R_j(\Ad_{\mathrm f})\cap J^\circ_j$, whose points are integral at the places prime to $q$, on which the multiplicity
$m_j$ (for $S=R_j(\Ad_{\mathrm f})$, $A=J^\circ_j$) divides a power of $q$, and on which $m_j\ne0$ at every point of bounded orbit for
$(\mathfrak q,j)$. Put $\Omega_0=J_M/J_0$ with $m_0=1$.
\lean{New.Seesaw.coverMultiplicity, New.Seesaw.coverMultiplicity_smul, New.Seesaw.coverMultiplicity_eq_index_stabilizer, New.Data.IsStageCut, New.Data.IsStageCut.integral, New.Data.IsStageCut.mem_iff_smul_mem, New.Data.IsStageCut.multiplicity_dvd_pow, New.Data.IsStageCut.ne_zero_of_hasBoundedOrbit, New.Data.cutPoints, New.Data.cutMultiplicity, New.Data.isStageCut_cutPoints, New.Data.mk_mem_cutPoints_zero}

\begin{proposition}[Integrals as weighted coset sums]\label{prop:coset-sum}
\begin{enumerate}[label=\textup{(\alph*)}]
\item Let $J\le A\le B$ be levels of $\U(V)(\Ad_{\mathrm f})$, $D$ compact with $D\subseteq B$ and, for each $x\in D$, $(xux^{-1})^{q^n}\in A$
for all $u\in J$ and some $n$; and let $W$ be a property such that each $x$ with $W(x)$ has $rx\in D$ for some $r\in S$. Then
there is a finite $T$, stable under $S\cap A$, with points in $B$, on which and at every point with $W$ the
multiplicity divides a power of $q$.
\item For a grading of $V$ and a level $A$ there is $\kappa$ such that: for every level $J$, cut $T$ stable under
$R_V(\Ad_{\mathrm f})\cap A$, and continuous, rationally invariant, integrable $f$ on $R_V(E^+)\backslash\U(V)(\Ad)$ with $m\ne0$ on its
support and $\int_{[R_V]}f(rt)\,dr=[J^\circ:J]\,a(t_\infty)b(t_f)$ for $t_f\in T$,
\[
  \int_{R_V(E^+)\backslash\U(V)(\Ad)}f=\kappa\sum_{t\in T}m(t)^{-1}\,b(t).
\]
\end{enumerate}
In particular cuts exist for every stage.
\end{proposition}
\lean{New.Seesaw.exists_finset_coverMultiplicity_dvd_pow, New.Grading.exists_integral_eq_sum_over_cut, New.Seesaw.exists_index_stabilizer_dvd_pow, New.Seesaw.exists_index_dvd_pow_of_pow_mem, New.Seesaw.integral_quotMeasure_eq_integral_withDensity, New.Seesaw.lintegral_indicator_inv_coverMultiplicity, New.Seesaw.integral_withDensity_mul_eq_mul, New.Seesaw.measure_union_cosets, New.Seesaw.lintegral_besselGroup_cutoff_mul, New.Seesaw.integral_indicator_smul_quotient, New.Data.exists_isStageCut, New.Data.exists_conj_pow_mem_stdLevel, New.Grading.exists_finite_representatives}

\begin{proof}
(a) \emph{Multiplicities as indices.} Let $C$ be a set of representatives, one in each $S$-orbit, of the cosets $xJ$ whose $S$-orbit meets $DJ/J$, and let $T$ be the union of the orbits of $C$ under $S\cap A$. The multiplicity is constant on $S$-orbits: left translation by $p^{-1}$ carries the classes counted at $pt$ to those counted at $t$. At
$t=xJ\in C$, the classes $r(S\cap A)$ with $r^{-1}t\in T$ are those with $r^{-1}t$ in the $(S\cap A)$-orbit of $t$, i.e.\ the orbit of the trivial class under the stabiliser $Z=S\cap xJx^{-1}$ of $t$; so the multiplicity is the index $[Z:Z\cap A]$. This index is finite, because $Z$ lies in the compact group $xAx^{-1}$ and $A$ is
open.

\emph{The index divides a power of $q$.} Every element of $Z$ is $xux^{-1}$ with $u\in J$, and by hypothesis $(xux^{-1})^{q^n}\in A$; so the $q^n$-th power of every element of $Z$ lies in $Z\cap A$. Let $N$ be the normal core of $Z\cap A$ in $Z$, a subgroup of finite index. In the finite group $Z/N$ every element has order dividing $q^n$, so
by Cauchy's theorem every prime dividing its order divides $q$, and its order, hence $[Z:Z\cap A]$, divides a power of $q$.

\emph{The cut.} $T$ is stable under $S\cap A$ by construction. It is finite, because it lies in the image of the compact set $AD$ in $H/J$ and $J$ is open; and its points lie in $B$, because $D\subseteq B$ and $A\subseteq B$. At a point of $T$ the multiplicity is that at a point of $C$, and if $W(x)$, then $rx\in D$ for some $r\in S$, so the
$S$-orbit of $xJ$ contains a point of $C$ and the multiplicity is again that at a point of $C$. In both cases it divides a power of $q$.

(b) \emph{A weight.} Let $\eta$ be the Haar measure on $S=R_V(\Ad_{\mathrm f})$ that gives $S\cap A$ the mass $1$. The function $\psi=\sum_{t\in T}m(t)^{-1}\mathbf1_{tJ}$ on $\U(V)(\Ad_{\mathrm f})$ satisfies
\[
  \int_S\psi(u^{-1}t)\,d\eta(u)=\begin{cases}1,&m(t)\ne0,\\0,&m(t)=0,\end{cases}
\]
because the multiplicity is constant on the orbit of $t$ and the integrand is the constant $m(t)^{-1}$ on the $m(t)$ classes of $S$ modulo $S\cap A$ that it counts and $0$ elsewhere. Let $\beta$ be a cutoff for the archimedean Bessel group, of integral one over each of its orbits, and $\alpha$ its archimedean part; then $w(g)=\alpha(g_\infty)\psi(g_f)$
has integral $1$ over each orbit of $R_V(\Ad)$ on which $m\ne0$, and $0$ over the others.

\emph{Unfolding.} The map $(t,r)\mapsto rt$ carries the measure $w(t)\,dt\times dr$ to the measure of $R_V(E^+)\backslash\U(V)(\Ad)$ on the set where $m\ne0$; as $f$ vanishes off that set,
\begin{align*}
  \int_{R_V(E^+)\backslash\U(V)(\Ad)}f&=\int_{\U(V)(\Ad)}w(t)\int_{[R_V]}f(rt)\,dr\,dt\\
  &=[J^\circ:J]\int_{\U(V)(\Ad)}\alpha(t_\infty)a(t_\infty)\,\psi(t_f)b(t_f)\,dt ,
\end{align*}
the second equality by the period hypothesis, as $\psi(t_f)=0$ unless $t_f\in T$. The integral of a product weight against a product function is the product of the two integrals. The finite one is $\sum_{t\in T}m(t)^{-1}b(t)\,\mathrm{vol}(J)$, and $\mathrm{vol}(J)=[J^\circ:J]^{-1}$ since $J^\circ$ has volume $1$; so the index cancels, and
$\kappa=\int\alpha a$ depends only on $a$, the grading and the level $A$ through $\eta$.

\emph{Existence of cuts.} Let $\mathfrak q$ be as in \cref{ss:support}. By the finiteness of the lattices with the bounds (\cref{ss:B-unfolding}) there is a finite set $C$ of points supported on $Q\cup\bar Q$ such that $R_j(\Ad_{\mathrm f})\,x\,J^\circ_j$ meets $C$ for every $x$ with bounded orbit for $(\mathfrak q,j)$. Apply (a) with $J=J_j$, $A=J^\circ_j$, $B$ the
group of points integral at the places prime to $q$, $D=CJ^\circ_j$, which is compact and contained in $B$, and $W$ the bounded orbit property. The power condition holds: for $x\in D$ and $u\in J_j$, at a place outside $Q\cup\bar Q$ the point $x$ is integral and $xux^{-1}$ is integral; at a place $w\in Q$, $u\equiv1$ modulo $w$, conjugation by $x$
keeps the elements close enough to $1$ integral, and high powers of elements $\equiv1$ are close to $1$ by lifting the exponent, so $(xux^{-1})^{q^n}$ is integral for some $n$. The set $T$ of (a) is then a cut: it is stable, its points are integral at the places prime to $q$, its multiplicities divide powers of $q$, and at a point of bounded
orbit the multiplicity divides a power of $q$, hence is not $0$.
\end{proof}

For a function $S$ of tuples of classes, the \emph{cut sum} is
\[
  \Sigma_kS=\sum_{c_0\in\Omega_0}\cdots\sum_{c_{k-1}\in\Omega_{k-1}}\ \prod_{j<k}m_j(c_j)^{-1}\,S(c_0,\dots,c_{k-1},1,1,\dots).
\]
It commutes with linear maps, and if $b\in\U(X_j)(\Ad_{\mathrm f})$ lies in the
level of $X_j$ its image in the Bessel group permutes $\Omega_j$ and preserves the multiplicities. If $Q=\emptyset$, each
$\Omega_k$ has at most $M^{fn_m^2}$ elements.
\lean{New.Data.cutSumAt, New.Data.cutSum, New.Data.map_cutSum, New.Data.cutSumAt_smul, New.Data.cutSum_comp_smul, New.Data.smul_mem_cutPoints_iff, New.Data.cutMultiplicity_smul, New.Data.inclXFin_gradingY_mem_besselFin, New.Data.card_le_residueMatrixCount, New.Data.mem_localLevel_of_mem_cutPoints, New.Data.exists_cutMultiplicity_dvd, New.Data.residueMatrixCount}

These properties are immediate. The cut sum is a finite sum with scalar coefficients. An element $b$ of the level of $X_j$ maps into $J_M\subseteq J^\circ_j$ (its reduction at the chosen places is $1$, and above $N$ the congruence on the last row holds, \cref{prop:graded-lattice}(d)) and into the Bessel group; for $j>0$ the cut is stable under
$R_j(\Ad_{\mathrm f})\cap J^\circ_j$ and the multiplicity is constant on $R_j(\Ad_{\mathrm f})$-orbits, and for $j=0$ left translation by an element of $J_M$ permutes $J_M/J_0$. If $Q=\emptyset$, the points over $\Omega_k$ are integral everywhere, and the principal level of $\pi_M$ lies in $J_k$; so $\Omega_k$ injects into the integral points
modulo $\pi_M$, a subgroup of the $n_k\times n_k$ matrices over $\OO_E/\pi_M$, which has $\Nm(\pi_M)^{n_k^2}=M^{fn_k^2}\le M^{fn_m^2}$ elements.

\subsection{The neighbour lattice}\label{ss:neighbour}
Let $P$ consist of one place of each conjugate pair above $N$, let $X$, $X'$, $X''$ be definite with last vectors $e$, $e''$
of $X\oplus\langle e\rangle$ and $X''\oplus\langle e''\rangle$, and let $\eta$ be the unit of norm $-1$. For residues $y$ modulo $N$ at $P$ the
\emph{neighbour translation} $n(y)$ is the matrix in $\Hom(X',X^-\oplus X'')$ with row $\eta\bar y/N$ at $e$, row $-\bar y/N$ at $e''$, and
zero elsewhere; it is supported at the places of $\bar P$. Scaling the rows of the standard lattice by factors $\kappa$ that
are $1$ or a norm-one adele $\delta$, trivial above $N$, and applying an explicit unitary gives the \emph{neighbour lattice}
$L'$. Away from $N$ it is the standard lattice with scaled rows; above $N$ it is a neighbour of the standard lattice, and the
translations $n(y)$ represent it modulo its integral points. A \emph{coset representative} is $n(y)+t$, with $t$ a
representative of a class of the elements of $L'$ that vanish above $N$.
\lean{New.Seesaw.PairIndex, New.Seesaw.PairIndex.lastSource, New.Seesaw.PairIndex.lastTarget, New.Seesaw.selfOrConj, New.Seesaw.vanishingAbove, New.Seesaw.mem_vanishingAbove, New.Seesaw.mulVec_mem_vanishingAbove, New.Seesaw.sub_idem_smul_mem_vanishingAbove, New.Seesaw.neighbourTranslation, New.Seesaw.neighbourLattice, New.Seesaw.latticeUnitary, New.Seesaw.coe_latticeUnitary, New.Seesaw.latticeUnitary_smul, New.Seesaw.latticeUnitary_smul_of_mem_vanishingAbove, New.Seesaw.neighbourUnitary, New.Seesaw.coe_neighbourUnitary, New.Seesaw.scalingUnitary, New.Seesaw.kroneckerOne, New.Seesaw.coe_kroneckerOne, New.Seesaw.kroneckerOne_smul, New.Seesaw.kroneckerOne_scalingUnitary_smul, New.Seesaw.neighbourMatrix, New.Seesaw.neighbourMatrix_apply, New.Seesaw.neighbourMatrix_apply_of_notMem, New.Seesaw.neighbourMatrix_mulVec, New.Seesaw.commute_neighbourMatrix_diagonal, New.Seesaw.neighbourAxis, New.Seesaw.neighbourAxis_mem, New.Seesaw.star_neighbourAxis_dotProduct, New.Seesaw.star_neighbourAxis_dotProduct_self, New.Seesaw.neighbourCoeff, New.Seesaw.neighbourCoeff_apply_of_notMem, New.Seesaw.mul_neighbourCoeff, New.Seesaw.neighbourCoeff_sub_lastSource, New.Seesaw.idem_conjClosure_mul_neighbourCoeff, New.Seesaw.star_normNegOne_mul_neighbourCoeff_sub_lastTarget, New.Seesaw.neighbourScale, New.Seesaw.neighbourScale_apply_of_notMem, New.Seesaw.star_mul_neighbourScale_mem, New.Seesaw.star_one_add_neighbourScale_mul, New.Seesaw.conjTranspose_mul_one_add_smul_vecMulVec, New.Seesaw.normNegOne_mem, New.Seesaw.IsScalingAtM.star_idem_mul_apply, New.Seesaw.IsScalingAtM.idem_mul, New.Seesaw.IsScalingAtM.apply_eq_one, New.Seesaw.IsScalingAtM.star_idem_mul, New.Seesaw.IsScalingAtM.star_mul_mem, New.Seesaw.IsScalingAtM.apply_mul_star_eq_one, New.Seesaw.IsScalingAtM.v_eq_one, New.Seesaw.IsScalingAtM.mul_star_self_eq_one, New.Seesaw.IsScalingAtM.eq_one_or_eq, New.Seesaw.IsScalingAtM.selfOrConj_mul_mem, New.Seesaw.IsScalingAtM.idem_mul_apply, New.Seesaw.IsSplitSquarefree.one_sub_idem_mul_inv_mem, New.Seesaw.IsSplitSquarefree.isConjDisjoint, New.Seesaw.IsSplitSquarefree.natCast_eq_mul_star, New.Seesaw.IsSplitSquarefree.mem_integralAdeles, New.Seesaw.IsSplitSquarefree.star_mem_integralAdeles, New.Seesaw.star_idem_mul_self, New.Seesaw.star_mul_star_inv, New.Seesaw.star_idem_mul_star_residueOverN, New.Seesaw.exists_class_star_idem_mul_mem_iff}

\begin{proposition}[Neighbour cosets]\label{prop:neighbour-cosets}
\begin{enumerate}[label=\textup{(\alph*)}]
\item The neighbour translations lie in $L'$, and every element of $L'$ is, modulo the integral lattice, $n(y)+t$ with unique
$y$ and $t$ vanishing above $N$; the elements of $L'$ vanishing above $N$ are the rows scaled by $\kappa$, and the number
$K$ of their classes divides a power of the norm of the denominator $\beta$ of $\kappa$.
\item A vector $(c,x)$ solving the gluing congruences (a \emph{neighbour vector}) has its entries at $e$, $e''$ congruent to the
coefficient $c$ modulo $\pi$, and the column of the neighbour matrix is such a vector.
\end{enumerate}
\end{proposition}
\lean{New.Seesaw.neighbourTranslation_mem_neighbourLattice, New.Seesaw.neighbourCosetRep, New.Seesaw.neighbourCosetRep_mem_neighbourLattice, New.Seesaw.existsUnique_sub_neighbourCosetRep_mem, New.Seesaw.existsUnique_add_neighbourTranslation_mem, New.Seesaw.exists_add_neighbourTranslation_mem_iff, New.Seesaw.mem_neighbourLattice_iff_of_mem_vanishingAbove, New.Seesaw.idem_mul_neighbour_add_translation_mem, New.Seesaw.star_idem_mul_neighbour_add_translation, New.Seesaw.idem_conjClosure_smul_neighbourTranslation, New.Seesaw.star_idem_smul_neighbourTranslation, New.Seesaw.awayLattice, New.Seesaw.AwayLatticeQuotient, New.Seesaw.awayLattice_le_weightedTranslations, New.Seesaw.card_awayLatticeQuotient_dvd, New.Seesaw.mem_weightedTranslations, New.Seesaw.conjTranspose_mulVec_mem_weightedTranslations, New.Seesaw.star_mem_weightedTranslations, New.Seesaw.IsNeighbourVector, New.Seesaw.IsNeighbourVector.mem_of_ne, New.Seesaw.IsNeighbourVector.star_mul_mem, New.Seesaw.IsNeighbourVector.mul_coeff, New.Seesaw.IsNeighbourVector.sub_target_mem, New.Seesaw.IsNeighbourVector.sub_source_mem, New.Seesaw.isNeighbourVector_col_neighbourMatrix}

\begin{proof}
\emph{The unitary.} Let $j$ be the adele that is $1$ on $P$ and $0$ elsewhere, $\eta$ a unit of $\hat\OO_E$ with $\bar\eta\eta=-1$, and $\rho=e+\bar\eta\bar je''$: it is $e+\bar\eta e''$ at the places of $\bar P$ and $e$ elsewhere. Since $j\bar j=0$, $\langle\rho,\rho\rangle=1$, and $\langle\rho,x\rangle=x_e+\eta jx_{e''}$. Let $A$ be the adele with $1+A=\pi$ on
$P$, $\bar\pi^{-1}$ on $\bar P$ and $1$ elsewhere; then $1+A$ has norm one, and $\bar\pi A$ is integral. The matrix $\gamma=1+A\rho\rho^*$ sends $x$ to $x+A\langle\rho,x\rangle\rho$; it multiplies $\rho$ by $1+A$ and fixes $\rho^\perp$, and as $\rho\rho^*$ is a hermitian idempotent, $\gamma^{\mathsf H}\gamma=1+\bigl((1+\bar A)(1+A)-1\bigr)\rho\rho^*=1$.
The neighbour lattice is the image of the standard lattice under $(\gamma\kappa)\otimes1$, which acts on each column; away from $P\cup\bar P$, $\gamma=1$, and above $N$ the factors $\kappa$ are $1$, so $\gamma$ commutes with the scaling.

\emph{Above $N$.} At a place of $P$, $\gamma$ is integral ($A=\pi-1$ and $\rho=e$) and $n(y)=0$. At a place of $\bar P$, $A=\bar\pi^{-1}-1$, $\langle\rho,x\rangle=x_e$, and $n(y)=(\eta\bar y/N)\rho$ because $\bar\eta\eta=-1$; so each column of $\gamma x+n(y)$ is
\[
  x+\rho\,\bigl(\bar\pi^{-1}x_e+\eta\bar y/N-x_e\bigr).
\]
Its entries off $e$, $e''$ are those of $x$; its entry at $e$ is $\bar\pi^{-1}(x_e+\eta\bar y/\pi)$, and its entry at $e''$ is $x_{e''}-\bar\eta x_e$ plus $\bar\eta$ times the entry at $e$. At a place of $\bar P$, $\pi$ is a unit and $\bar\pi$ a uniformiser, so this column is integral exactly when $x$ is integral and $x_e\equiv-\eta\bar y/\pi$ modulo $\bar\pi$; if
$x$ is integral there, exactly one residue $y$ modulo $N$ at the conjugate place of $P$ satisfies this, and if not, none does. Hence: (i) every $n(y)$ lies in $L'$, since with $x=-\gamma^{-1}n(y)$ the column $\gamma x+n(y)=0$ is integral, so $x$ is integral at $\bar P$, and $n(y)$, being supported at $\bar P$, is the image of the part of $-x$ at
$\bar P$; and (ii) for every element $\gamma x_0$ of $L'$ ($x_0$ integral, $\kappa=1$ above $N$) there is exactly one $y$ for which $\gamma x_0+n(y)$ is integral at the places above $N$.

\emph{Away from $N$.} A matrix that vanishes above $N$ lies in $L'$ exactly when $\kappa^{-1}x=\bar\kappa x$ is integral, i.e.\ these are the rows scaled by $\kappa$; such a matrix is $\kappa$ times an integral matrix, so its rows times $\beta$ in the first block and $\bar\beta$ in the second are integral. Reducing the conjugate of $\bar\kappa x$ modulo $\beta$
defines an injective homomorphism from these elements modulo the integral ones into the matrices over $\hat\OO/\beta$: if the reduction is $0$, then $x$ is $\bar\beta\kappa$ times an integral matrix, hence integral. Hence $K$ divides the number of such matrices, a power of $\Nm(\beta)$.

\emph{Coset representatives.} Let $z\in L'$. Its part above $N$ decides $y$ by (ii), and $z+n(y)$ is integral above $N$; its part away from $N$ is an element of $L'$ vanishing above $N$, whose class modulo the integral matrices decides $t$. Hence $z\equiv n(-y)+t$ with unique $(y,t)$; and $n(y)+t\in L'$ for all $y$ and $t$ by (i).

(b) Let $x$ be the column of $\gamma$ at an index $m$ and $c$ its coefficient: $A\pi^{-1}\bar j$ for $m=e$, $A\pi^{-1}\eta j$ for $m=e''$ and $0$ otherwise. The coefficient is supported above $N$, where $\delta=1$, so $\delta c=c$. The entries of $x$ off $e$, $e''$ are those of the identity; all of $\bar\pi x$ is integral as $\bar\pi A$ is. A direct computation
gives $c-\pi^{-1}x_e=-\varpi$ for $m=e$ and $0$ otherwise, and $\bar\eta c-\pi^{-1}x_{e''}=-\varpi$ for $m=e''$ and $0$ otherwise, with $\varpi=j+(1-j)\pi^{-1}$, which is integral because $\pi$ is a unit off $P$. These are the gluing congruences: the entries at $e$ and $e''$, divided by $\pi$, are congruent to $c$ and $\bar\eta c$.
\end{proof}

\begin{theorem}[Finite evaluation of the vacuum]\label{thm:finite-evaluation}
Let $Y\to Y'\to Y''$ be a pair of edges with $N=\pi\bar\pi$ split and squarefree, let the point of $Y'$ span its graded
lattice, let every neighbour vector lift to the lattice of $V^-\oplus V''$, and let the denominators match through $\kappa$.
\begin{enumerate}[label=\textup{(\alph*)}]
\item The image in $I^\perp/I$ of $L\cap I^\perp$, for the standard lattice $L$ of $\Hom(V',V^-\oplus V'')$ and the isotropic subspace
$I$ of the splitting, is the neighbour lattice; and $Z_f(\mathrm{vac})=c\sum_{(y,t)}\varrho(n(y)+t)\,\mathrm{vac}$ with $c\ne0$.
\item There is $C$ such that, for $n$ prime to $N$ and translations $x$ with denominators $\beta$ on $V^-$ and $\bar\beta$ on $V''$ away from
$n$, the finite evaluation of $\varrho(x)\mathrm{vac}$ is a finite combination of $C\cdot\mathrm{Glue}(\varrho(t_1)\mathrm{vac}\otimes\varrho(t_2)\mathrm{vac})$ with
$\bar\beta t_1$, $\bar\beta t_2$ integral away from $n$, whose coefficients have a denominator away from $n$ depending only on $\beta$
and the index sets, and absolute sum at most $1$ under every automorphism of $\C$.
\item The hypotheses hold for the middle and the last pairs of the chain $Y_0\to\dots\to Y_m$, with $\kappa$ the last scaling
and $\beta=\bar\pi_M$.
\end{enumerate}
\end{theorem}
\lean{New.Seesaw.DoubleLiftData.zeroTermLattice_splitting_eq, New.Seesaw.DoubleLiftData.exists_zeroTermFin_splitting_vacElem_eq_sum, New.Seesaw.isAdmissibleEval_finEval, New.Seesaw.DoubleLiftData.exists_finEval_finHeisenbergOp_vacElem, New.Seesaw.DoubleLiftData.LiftsNeighbours, New.Seesaw.DoubleLiftData.MatchesDenominators, New.Seesaw.DoubleLiftData.latticeLift, New.Seesaw.DoubleLiftData.isIntMatrix_latticeLift, New.Seesaw.DoubleLiftData.latticeLift_eq_latticeUnitary_add_isotropic, New.Seesaw.DoubleLiftData.conjTranspose_latticeLift_mulVec_stackedStringMapFin, New.Seesaw.DoubleLiftData.conjTranspose_latticeLift_mulVec_mem, New.Seesaw.DoubleLiftData.latticeUnitary_mulVec_conjTranspose_latticeLift_mulVec, New.Seesaw.DoubleLiftData.liftVector, New.Seesaw.DoubleLiftData.neighbourLift, New.Seesaw.DoubleLiftData.planeCoeff, New.Seesaw.DoubleLiftData.isotropicCoeff, New.Seesaw.DoubleLiftData.evalAt_liftVector_eq_zero, New.Seesaw.DoubleLiftData.evalAt_neighbourLift_eq_zero, New.Seesaw.DoubleLiftData.conjTranspose_iFin_mulVec_mem, New.Grading.gradeZeroRows, New.Grading.isIntMatrix_gradeZeroRows, New.Grading.conjTranspose_gradeZeroRows, New.Seesaw.inv_xFin_mulVec_mem_of_spansLattice, New.Seesaw.flipScale, New.Seesaw.coe_flipScale, New.Seesaw.inv_xFin_flipSum, New.Seesaw.inv_xFin_flipSum_mulVec_mem_iff, New.Seesaw.lastScaling, New.Seesaw.liftsNeighbours_middleDoubleLift, New.Seesaw.matchesDenominators_middleDoubleLift, New.Seesaw.liftsNeighbours_lastDoubleLift, New.Seesaw.matchesDenominators_lastDoubleLift, New.Seesaw.zeroTermLattice, New.Seesaw.mem_zeroTermLattice_iff, New.Seesaw.mem_zeroTermLattice_iff_selfDualAddChar, New.Seesaw.mem_map_stdLattice_iff_selfDualAddChar, New.Seesaw.map_stdLattice_le_zeroTermLattice, New.Seesaw.zeroTermLattice_eq_map_of_le, New.Seesaw.exists_zeroTermLattice_eq_map, New.Seesaw.relIndex_map_stdLattice_ne_zero, New.Seesaw.relIndex_zeroTermLattice_ne_zero, New.Seesaw.exists_zeroTermFin_vacElem_eq_sum, New.Seesaw.exists_zeroTermFin_finHeisenbergOp_eq_smul, New.Seesaw.exists_zeroTermFin_finHeisenbergOp_add, New.Seesaw.exists_finHeisenbergOp_add_vacElem, New.Seesaw.finHeisenbergOp_zeroTermFin_vacElem, New.Seesaw.zeroTermFin_finHeisenbergOp_finFrameVector, New.Seesaw.exists_sub_orthogonal_of_zeroTermFin_ne_zero, New.Seesaw.exists_eq_add_of_zeroTermFin_ne_zero, New.Seesaw.finFrameVector_zero, New.Seesaw.conjTranspose_isotropicBasisFin_iFin_mulVec_finFrameVector, New.Seesaw.exists_eq_iFin_add_isotropicBasisFin_of_mulVec_eq_zero, New.Seesaw.conjTranspose_mulVec_iFin_add_isotropicBasisFin, New.Seesaw.exists_isIntMatrix_iFin_mul_add_isotropicBasisFin_mul, New.Seesaw.exists_omega_vacElem_eq_sum, New.Seesaw.ratCoordMatrix_sum, New.Seesaw.gluedTranslate, New.Seesaw.stdPairSpace, New.Seesaw.IsAdmissibleEval, New.Seesaw.isAdmissibleCombination_finHeisenbergOp_sum_neighbourCosetRep, New.Seesaw.exists_sum_finHeisenbergOp_neighbourTranslation_eq, New.Seesaw.finHeisenbergOp_comm_neighbourTranslation, New.Seesaw.finHeisenbergOp_neighbourTranslation_add, New.Seesaw.star_neighbourTranslation_dotProduct_eq_zero, New.Seesaw.arrowTensor_doubleLiftGluing, New.Seesaw.arrowTensor_tmul, New.Seesaw.finHeisenbergOp_vacElem_eq_smul_arrowTensor, New.Seesaw.exists_finHeisenbergOp_vacElem_eq_of_coprime, New.Seesaw.v_selfOrConj_eq_one, New.Seesaw.semichar_eq_one_of_smul_eq}

\begin{proof}
(a) \emph{The zero term of the vacuum.} Let $L$ be the standard lattice of $\Hom(V',V^-\oplus V'')$, $I$ the isotropic subspace of the splitting, with finite frame $k$, and $i$ the finite embedding of $\Hom(X',X^-\oplus X'')$. The lattice $L^\circ$ of the zero term is the set of $w$ with $iw+kb\in L$ for some $b$, the image of
$L\cap I^\perp$ in $I^\perp/I$. There are a finite unitary point $g$ and a matrix $c$ with $ig+kc$ integral (take a point adapted to the splitting); for such $g$ and $x$ integral, $(ig+kc)x$ is a lattice vector of class $gx$, and a vector $w$ lies in $L^\circ$ if and only if $g^{-1}w$ is integral. Hence $L^\circ$ is the image of the standard lattice under a
unitary point, hence self-dual. The finite zero term of the vacuum is fixed by the translations by $L^\circ$ and is not zero, so it is a nonzero multiple of $\omega(g)\mathrm{vac}$, that is, of the sum of the translates of the vacuum by representatives of $L^\circ$ modulo its integral points.

\emph{Identification with the neighbour lattice.} Let $B$ be the matrix of $x\mapsto A_\nu xR$, where the columns of $A_\nu$ are the lifts to $V^-\oplus V''$ of the columns of the neighbour matrix with their coefficients, and $R$ consists of the rows of the point of $V'$ at $X'\oplus\langle e'\rangle$. The columns of $A_\nu$ lie in the lattice by the
hypothesis on neighbour vectors and \cref{prop:neighbour-cosets}(b), and $R$ is integral because the string of $Y'$ has even length, so that the graded basis is the identity on $X'\oplus\langle e'\rangle$; hence $B$ is integral. $B^{\mathsf H}$ first restricts a map to $X'\oplus\langle e'\rangle$, where the stacked map vanishes, so $B^{\mathsf H}$
kills the stacked map. In the part of grade zero a lifted column is the column of $\kappa\gamma$ plus an isotropic vector of the planes, so $Bx=i(\gamma\kappa x)+k(\cdot)$: $Bx$ lies in $L\cap I^\perp$ with class $(\gamma\kappa\otimes1)x$. Hence the neighbour lattice is contained in $L^\circ$. Both are images of the standard lattice under unitary
points, so both are self-dual, and a self-dual lattice contained in another one equals it ($L_1\subseteq L_2$ gives $L_2=L_2^\vee\subseteq L_1^\vee=L_1$). Hence $L^\circ$ is the neighbour lattice, and by \cref{prop:neighbour-cosets}(a) its representatives modulo integral points are the $n(y)+t$.

(b) \emph{Reduction away from $N$.} The places above $N$ do not divide $n$, and $\beta$ is a unit there, so $x$ is integral above $N$; its part $l$ above $N$ is a lattice vector, and $\varrho(x)\mathrm{vac}$ is a root of unity times $\varrho(x-l)\mathrm{vac}$. Hence we may assume that $x$ vanishes above $N$.

\emph{One translate.} The evaluation translates by the stacked map $\tau$ and then takes the finite zero term. The translations by $\tau$ and $x$ compose to the translation by $\tau+x$ up to a root of unity. If $Z_f(\varrho(\tau+x)\mathrm{vac})\ne0$, then $\tau+x=iw+kb+l$ with $l$ integral, and $Z_f(\varrho(\tau+x)\mathrm{vac})$ is a root of unity times
$\varrho(w)Z_f(\mathrm{vac})$; the class $w$ differs from $gB^{\mathsf H}(\tau+x)$ by an element of $L^\circ$, so one may take $w=gB^{\mathsf H}(\tau+x)=gB^{\mathsf H}x$, as $B^{\mathsf H}\tau=0$. Here $g=(\gamma\kappa)\otimes1$, and since the isotropic part of the lift is supported above $N$, where $x$ vanishes, $gB^{\mathsf H}x=i^{\mathsf H}x$
is the orthogonal projection of $x$. The projection keeps the denominators: $B^{\mathsf H}$ maps a vector with row denominators $\beta$, $\bar\beta$ away from $n$ to one with denominators $\beta\kappa$, $\bar\beta\kappa$, because the denominators match through $\kappa$ (an entry of $B$ joining indices with different valuations of the denominators
vanishes at the conjugate place), and $g$ multiplies back by $\kappa$. Hence the evaluation of $\varrho(x)\mathrm{vac}$ is $0$ or a root of unity times the left inverse of the join applied to $c\,\varrho(x')\sum_{(y,t)}\varrho(n(y)+t)\mathrm{vac}$, with $x'=i^{\mathsf H}x$ vanishing above $N$ and of row denominators $\beta$, $\bar\beta$ away
from $n$.

\emph{The sum over the neighbour lattice.} The translation $n(y)$ is supported at $\bar P$ and is orthogonal to every vector vanishing above $N$, so it commutes with the translations by such vectors, and the translation by $n(y)+v$ is the product of the two translations. For each class $t$, $\varrho(x')\varrho(t)\mathrm{vac}$ is a root of unity times
$\varrho(x'+t)\mathrm{vac}$, and the translate of the vacuum by a matrix $z$ is a root of unity times the join of $\varrho(t_1)\mathrm{vac}\otimes\varrho(t_2)\mathrm{vac}$, with $t_1$ the unit $\eta$ times the adjoint of the first block of $z$ and $t_2$ its second block; as the first block of $x'+t$ has denominator $\beta$ and the second $\bar\beta$ away
from $n$, $\bar\beta t_1$ and $\bar\beta t_2$ are integral away from $n$. Under the join, $\sum_y\varrho(n(y))$ is the gluing operator of \cref{prop:glued-datum}: the adjoint times $\eta$ turns the row $\eta\bar y/N$ into the last column $-y/N$ of the first edge, and the row $-\bar y/N$ is the last row of the second. Hence the evaluation is $c$ times a
root of unity times $\sum_t\zeta_t\,\mathrm{Glue}(\varrho(t_1)\mathrm{vac}\otimes\varrho(t_2)\mathrm{vac})$, a sum of $K$ terms with roots of unity $\zeta_t$. Writing it as $\sum_t(\zeta_t/K)\cdot K\,\mathrm{Glue}(\cdots)$ and $C=cK$, the coefficients have absolute sum $1$ under every automorphism, and their denominator divides the number of integral
matrices modulo $\beta$ of the shape of $\Hom(X',X^-\oplus X'')$, because $K$ does; this depends only on $\beta$ and the index sets.

(c) \emph{Middle pairs.} The lift of a neighbour vector $(c,x)$ has the component $sx$ on $X\oplus\langle e\rangle$, with $s$ the idele by which negating the form scales the point, the component $x$ on $X''\oplus\langle e''\rangle$, and the coefficient $suc$ at the isotropic vector $z+z''$ formed by the middle vectors. A vector of $V$ with component
$x$ on $X\oplus\langle e\rangle$, coordinate $u(\pi^{-1}x_e+c')$ at the middle vector and $0$ at the other string vectors lies in the graded lattice as soon as $c'$, $\bar\pi x_e$ and the other entries of $x$ are integral, because its coordinates in the graded basis are $x$ off $e$, $\bar\pi x_e$ at $v$ and $c'$ at the middle vector. The
congruence $x_e/\pi\equiv c$ makes $c'=c-\pi^{-1}x_e$ integral. On $V''$ the same test applies with $u''$ and $\bar\eta c$, because $su_{2i}=\bar\eta\delta u_{2i+2}$: the scale of $Y_{2i}$ is that of $Y_{2i+2}$ times $a/\bar a$ (\cref{prop:lattice-Y}), and $s\,a/\bar a=\bar\eta\,c/\bar c=\bar\eta\delta$ for the chosen idele
$c$, the factors $\alpha$, $\bar\alpha$ cancelling. The inclusion of $X^-\oplus X''$ keeps each index in its block, so the denominators match with $\kappa=1$.

\emph{The last pair.} Here no plane is removed, and $e''$ is $z/p$ for the string vector $z$ of $V$, with $N=p\bar p$. The lift of $(c,x)$ has the component $sx$ on $X\oplus\langle e\rangle$, $x$ on $X''$ and $\delta x_{e''}/p$ at $z$; in the membership test the terms in $c$ cancel because $\bar\eta\eta=-1$, and the remaining entries are integral by
the congruences. The only index that changes its block is $e''$, whose denominator $\bar\beta$ is scaled by $\delta$ in the last scaling, and $\bar\beta\delta$ has the valuations of $\beta$ since $\delta=c/\bar c$ for an idele $c$ with the valuations of $\beta$. Finally $\delta$ has norm one and is $1$ above $N$, where $\beta$ is a unit, and the
factors $1$, $\delta$ have the required denominators, so these are scalings.
\end{proof}

\subsection{Admissible combinations}\label{ss:admissible}
For a family of vectors $v(x)$ indexed by a set $\mathcal X$, a vector $z$ is an \emph{admissible combination} with data
$(A,n,d,p,B)$ if $z=\sum_{x\in A}r(x)v(x)$ with $r$ finitely supported in $A$, $d\,r(x)$ integral away from $n$ for all $x$, and,
if the proposition $p$ holds, $\sum_x\abs{\sigma r(x)}\le B$ for every automorphism $\sigma$ of $\C$. The \emph{translations}
with denominator $d$ away from $q$ are the finite adelic vectors $x$ with $dx$ integral at the places prime to $q$.
\lean{New.Seesaw.IsAdmissibleCombination, New.Seesaw.integralAway, New.Seesaw.div_mem_integralAway, New.Seesaw.IsAdmissibleCoeffs, New.Seesaw.IsAdmissibleCoeffs.support_subset, New.Seesaw.IsAdmissibleCoeffs.sum_norm_le, New.Seesaw.IsAdmissibleCoeffs.integral, New.Seesaw.translations, New.Seesaw.weightedTranslations, New.Seesaw.star_idem_smul_mem_translations, New.Seesaw.mulVec_mem_translations, New.Seesaw.mem_translations_of_mem_stdLattice, New.Seesaw.localLevel_le_integralAbove_of_coprime}

\begin{proposition}[Rules for admissible combinations]\label{prop:admissible-rules}
Admissible combinations obey the following rules.
\begin{enumerate}[label=\textup{(\alph*)}]
\item \emph{Substitution.} Let $z$ be an admissible combination of the vectors $v(x)$ with data $(A,n,d,p,B)$, and let each
$v(x)$, $x\in A$, be one of the vectors $v'(x')$ with data $(A',n,d',p,B')$. Then $z$ is an admissible combination of the
$v'(x')$ with data $(A',n,dd',p,BB')$. Scalars integral away from $n$, of
absolute value at most $1$ under all $\sigma$ when $p$ holds, linear and bilinear maps, pure tensors and finite sums obey
the corresponding rules; a root of unity may be absorbed, and a nonzero scalar may be normalised away.
\item \emph{Moved vacua.} Let $g$ be a finite unitary point with multiplier $c\,C$ and $A$ contain all $g(x+z)$, $z$ integral. If
$d\,c/\abs{\Delta_f(g)}^2$ is integral away from $q$ and $\abs{\sigma c}\le1$ when $p$ holds, then $\omega(1,g)\varrho(x)\mathrm{vac}$ is an
admissible combination of the $C\varrho(y)\mathrm{vac}$, $y\in A$, with $(A,q,d,p,1)$. In particular this holds for $g$ integral away from $q$.
\item \emph{Functions modulo an ideal.} The datum of a function modulo $\mathfrak m$ (\cref{prop:mod-datum}) is an admissible
combination of translates of the vacuum, with denominator the number of residue matrices.
\end{enumerate}
\end{proposition}
\lean{New.Seesaw.IsAdmissibleCombination.bind, New.Seesaw.IsAdmissibleCombination.bind_one, New.Seesaw.IsAdmissibleCombination.of_sum, New.Seesaw.IsAdmissibleCombination.smul_single, New.Seesaw.IsAdmissibleCombination.single, New.Seesaw.IsAdmissibleCombination.map, New.Seesaw.IsAdmissibleCombination.smul, New.Seesaw.IsAdmissibleCombination.smul_of_eq_zero_or_isOfFinOrder, New.Seesaw.IsAdmissibleCombination.of_eq_isOfFinOrder_smul, New.Seesaw.IsAdmissibleCombination.map_bilinear, New.Seesaw.IsAdmissibleCombination.smul_mem, New.Seesaw.IsAdmissibleCombination.normUnit_smul, New.Seesaw.IsAdmissibleCombination.pure, New.Seesaw.isAdmissibleCombination_weil_vacElem, New.Seesaw.isAdmissibleCombination_weil_of_mem_localLevel, New.Parameters.norm_finMultiplier_eq_one, New.AutChar.HasWeight.norm_apply_detFin_eq_one, New.Seesaw.card_residueRing_pi_dvd, New.Seesaw.isAdmissibleCombination_schrodingerOp_of_reduction, New.Seesaw.isAdmissibleCombination_schrodingerOp, New.Seesaw.exists_arrowProduct_eq_smul_finHeisenbergOp}

\begin{proof}
(a) Write $z=\sum_xr(x)v(x)$ and $v(x)=\sum_{x'}r'_x(x')v'(x')$. Then
\[
  z=\sum_{x'}\Bigl(\sum_xr(x)r'_x(x')\Bigr)v'(x'),
\]
a finite sum over $A'$. The ring of complex numbers integral away from $n$ (those whose images under all homomorphisms $\C\to\C_l$, $l\nmid n$, are integral) contains $dr(x)$ and $d'r'_x(x')$, hence
$dd'\sum_xr(x)r'_x(x')$. If $p$ holds, $\sum_{x'}\abs{\sigma\sum_xr(x)r'_x(x')}\le\sum_x\abs{\sigma r(x)}\sum_{x'}\abs{\sigma r'_x(x')}\le BB'$. A linear map sends $\sum r(x)v(x)$ to $\sum r(x)L(v(x))$ with the same coefficients. A scalar $c$ as stated multiplies the coefficients by $c$, and a single term $cv(x)$ is a combination of size $\abs{\sigma c}$. For a
bilinear map, expand $z$ in the first variable and $z'$ in the second: the data are $(A\times A',dd',BB')$; pure tensors follow by induction on the number of factors, the empty tensor being $1$. A root of unity $\zeta$ is integral and $\abs{\sigma\zeta}=1$. If $d\ne0$ and $z$ is a combination of the $Cv(x)$ with data $(A,n,d,p,B)$, then
$(d/(aC))\,az=\sum_xdr(x)v(x)$ has denominator $1$ and size $dB$, which is the normalisation (and $0$ if $aC=0$).

(b) Let $\Lambda=\hat\OO^n$ and $H=\Lambda\cap g^{-1}\Lambda$, of index $\abs{\Delta_f(g)}^2$ in $\Lambda$. The Weil operator of $g$ sends the vacuum to $\Delta_f(g)\abs{\Delta_f(g)}^{-2}\sum_{z\in\Lambda/H}\varrho(gz)\mathrm{vac}$ (\cref{lem:galois-law}), and $\omega(1,g)\varrho(x)=\varrho(gx)\omega(1,g)$; the translations by $gx$ and
$gz$ compose to the translation by $g(x+z)$ up to a root of unity $\zeta_z$, a value of the standard character on a rational multiple. Hence
\[
  \omega(1,g)\varrho(x)\mathrm{vac}=\sum_{z\in\Lambda/H}\zeta_z\,\frac c{\abs{\Delta_f(g)}^2}\,C\varrho(g(x+z))\mathrm{vac},
\]
with $\abs{\Delta_f(g)}^2$ terms; $d$ times each coefficient is integral away from $q$ by hypothesis, and as $\abs{\Delta_f(g)}^2$ is a natural number, the absolute sum of the $\sigma$-conjugates is $\abs{\sigma c}\le1$ when $p$ holds. If $g$ is integral at the places prime to $q$, take $C=1$, $c=\Delta_f(g)$:
$\Delta_f(g)=\upsilon(\det u)\abs{\det u}^{-1/2}$ for a matrix $u$ with $u\hat\OO^n=\hat\OO^n\cap g^{-1}\hat\OO^n$, and above a prime $l\nmid q$ the point $g$ and its inverse are integral, so $u$ may be taken $1$ there and $\det u$ is a unit above $l$; hence $\Delta_f(g)$ and its conjugates are units away from $q$, and
$c/\abs{\Delta_f(g)}^2=1/\overline{\Delta_f(g)}$ is integral away from $q$. If $q=1$, then $g$ is integral everywhere and $\Delta_f(g)=1$.

(c) Expand the function $\Psi$ modulo $\mathfrak m$ in indicators of points. The datum of the indicator of $y_0$ is the mean, over the vectors modulo $\mathfrak m$, of roots of unity times the translates of the integral vacuum by the vectors $\bar g^{-1}$ times integral vectors, supported at the places of $\mathfrak m$ (finite Fourier inversion, \cref{ss:B-weil},
(W6)). Hence the coefficients are $\Psi(y_0)$ times roots of unity divided by the number of vectors modulo $\mathfrak m$, whose absolute sums are at most $\sum\abs{\sigma\Psi}$, and multiplied by that number they are integral away from $q$.
\end{proof}

\begin{theorem}[The evaluated data are glued vacua]\label{thm:good-datum}
There is $C$ such that for every tame character $\vartheta$ and all points $y_k$ of $\U(Y_k)(\Ad_{\mathrm f})$ integral at the places prime
to $q$, the pure tensor $\phi_0(y)\otimes\dots\otimes\phi_\varrho(y)$ is an admissible combination of the vectors $C\,\mathrm{Glue}(x)$, where
$x$ runs over the tuples of translations with denominator $\bar\pi_M$ away from $q$ (the admissible translations), with
denominator $(M^{fn_m^2})^{2m}$ away from $q$ and, if $Q=\emptyset$, size $1$.
\end{theorem}
\lean{New.Data.exists_isAdmissibleCombination_gluedVacuum, New.Data.exists_isAdmissibleCombination_evaluatedDatum, New.Data.isAdmissibleCombination_pointDatumY, New.Data.isAdmissibleCombination_translatedPointDatumY, New.Data.isAdmissibleCombination_arrowProduct_translatedPointDatumY, New.Data.isAdmissibleEval_doubleLiftEvalY, New.Data.absNorm_pow_mul_card_dvd_residueMatrixCount_pow, New.Data.piM_mul_normQ_mem_ideal, New.Data.chosenRatio, New.Data.chosenRatio_mul_star, New.Data.chosenRatio_apply_of_N_mem, New.Data.v_piM_eq_one_of_N_mem, New.Data.v_star_piM_mul_chosenRatio, New.Data.flipScale_mul_divStar_ideleA, New.Data.isScalingAtM_of_eq_one_or_eq, New.Data.isScalingAtM_one, New.Data.isScalingAtM_lastScaling}

\begin{proof}
\emph{The data.} Let $\mathfrak m$ be the modulus of the twist, the product of the chosen places and the places of $Q$; it divides $\pi_Mq$. The datum $\Upsilon_{k+1}$ is the Schr\"odinger datum of the indicator of the class of the inclusion matrix modulo $\mathfrak m$, applied to the vacuum; by
\cref{prop:admissible-rules}(c), with the factor $\pi_M$ in the role of $g$, it is a combination of translates $\varrho(t)\mathrm{vac}$ with $\bar\pi_Mt$ integral away from $q$, whose coefficients times the number of matrices modulo $\pi_M$ are integral away from $q$, and whose size is $1$ (the indicator of one class has $\sum\abs{\sigma\Psi}=1$),
whatever the proposition $p$.

\emph{Moving.} The pair $(y_k,y_{k+1})$ acts on $\Hom(Y_k,Y_{k+1})$ as a scalar $c$ times the Weil operator of one finite point $g$, integral away from $q$ as the $y$ are. The scalar is a quotient of values of the characters of the lift at the determinants of $y_{k+1}$ and $y_k$; these characters have a weight, so at a finite point
integral above $l\nmid q$ their values are $l$-adic units, and if $Q=\emptyset$ the points are integral everywhere and the values are products of values of the tame parts, roots of unity, of absolute value $1$ under every automorphism. Expanding the datum in translates and moving each by $g$ (\cref{prop:admissible-rules}(b)), the translated datum
is an admissible combination of translates of the vacuum with the same denominators and, if $Q=\emptyset$, size $1$.

\emph{The join.} The join is bilinear, and the join of $\varrho(t_1)\mathrm{vac}$ and $\varrho(t_2)\mathrm{vac}$ is a root of unity times $\varrho(z)\mathrm{vac}$, where the first block of $z$ is minus the unit of norm $-1$ times the adjoint of $t_1$ and the second block is $t_2$; if $\bar\pi_Mt_1$ and $\bar\pi_Mt_2$ are integral away from $q$, the rows of $z$
have denominator $\pi_M$ on the block of $Y_{2i}$ and $\bar\pi_M$ on that of $Y_{2i+2}$. Hence by the bilinear rule the join of the two translated data of the pair $i$ is a combination of translates $\varrho(z)\mathrm{vac}$ with these denominators, whose denominator is the product of the numbers of matrices modulo $\pi_M$ of the two shapes, and of
size $1$ if $Q=\emptyset$.

\emph{The evaluation.} No place of $Q$ lies above $N$, so $q$ is prime to $N$, and \cref{thm:finite-evaluation}(b),(c) applies with $n=q$ and $\beta=\bar\pi_M$: each translate of the joined vacuum goes to a combination of the vectors $C\,\mathrm{Glue}(t_1,t_2)$ with $\bar\beta t_1$, $\bar\beta t_2$ integral away from $q$, with denominator the number of
integral matrices modulo $\beta$ of the shape of $\Hom(X_{2i+1},X_{2i}^-\oplus X_{2i+2})$ and size $1$. By substitution, $\phi_i(y)$ is a combination of the $C\,\mathrm{Glue}(t_1,t_2)$ whose denominator is the product of the three numbers of residue matrices; all three are powers of $\Nm(\pi_M)=M^f$, with exponents the numbers of entries, which add up to at
most $4n_m^2$ (the third shape has at most $2n_m^2$ entries, as all ranks are at most $n_m$), so the product divides $(M^{fn_m^2})^4$.

\emph{The tensor.} By the rule for pure tensors, $\phi_0(y)\otimes\dots\otimes\phi_\varrho(y)$ is a combination of the tensors of glued translates with denominator $(M^{fn_m^2})^{4(\varrho+1)}=(M^{fn_m^2})^{2m}$ and, if $Q=\emptyset$, size $1$. The tensor of the glued translates of the vacua is a fixed constant (a product of the factors relating the integral vacuum
of each edge to its vacuum) times $\mathrm{Glue}(x)$ for the tuple $x$ of translations, which gives the constant $C$ of the statement.
\end{proof}

\subsection{The seesaw identity}\label{ss:seesaw-proof}
For $j<m$ and a conjugate form $z$ on $Y_{j+1}$, the \emph{descent} of $z$ is $[J^\circ_j:J_j]$ times the theta integral
$g\mapsto\int_{[\U(Y_{j+1})]}\theta(\omega(g,h)(\varphi_{\rm harm}\otimes\Upsilon_{j+1}))\,z(h)\,dh$; it has level $J_j$, because the datum is fixed by
$J_j$, and lift and descent are transposes: $\int z\cdot L_j f=\int(\text{descent }z)\cdot f$. Conjugate forms on $Y_{j+1}$, $j<m$, are
rapidly decreasing. At the top, restriction to $Y_m$ is read in automorphic forms.
\lean{New.Data.ConjFormY, New.Data.transposedLiftY, New.Data.transposedLiftY_apply, New.Data.integral_mul_liftY, New.Data.twistedPairWeil_datumY_of_mem_levelY, New.Data.hasLevel_transposedLiftY, New.Data.isRapid_conjFormY, New.Data.restrictToFormY, New.Data.lowerChain_eq_of_lowerChainUpTo_eq, New.Data.unfoldedPeriod, New.Data.integrable_unfoldStage_mul}

\begin{proposition}[The outer integrals]\label{prop:outer-integrals}
\begin{enumerate}[label=\textup{(\alph*)}]
\item For $k\le2\varrho+1$, a grading of $Y_{k+1}$ and a conjugate form $z$ on $Y_{k+1}$, the period
$\Xi_{k+1}=\int_{R_{k+1}(E^+)\backslash\U(Y_{k+1})(\Ad)}\mathcal W_{k+1}z$ converges, and
$\int_{[R_k]}\mathcal W_k(rg_k)\,(\text{descent }z)(rg_k)\,dr=[J^\circ_k:J_k]\,\Xi_{k+1}$.
\item $\int_{[\U(Y_1)]}z\cdot L_0(\mathcal P\varphi)=[J^\circ_0:J_M]\sum_{c\in\Omega_0}\Xi_1(g,c)$.
\item For every function $a$ of the archimedean outer variables there is $c$ such that, whenever $\Xi_{k+1}=a(\xi)S(y)$ for all $y$ with
$y_0\in J_M$ and $y_m=1$, with $S$ a function of the classes of $y$ and $\varphi$ ramified,
\[
  \int_{[\U(Y_{k+1})]}z\cdot L_k\circ\dots\circ L_0(\mathcal P\varphi)=c\,\Sigma_{k+1}S .
\]
\item For a linear form $t$ on the values of forms on $X_m$ there is $C$ such that, if $F(\iota k)=t(\psi(\mathrm{diag}(k,1)))$ on $\U(Y_m)$ for a form
$\psi$ on $X_m$, then $\int_{[\U(Y_m)]}F\bar x=C\langle x,\mathrm{res}\,\psi\rangle$ for every form $x$ on $Y_m$.
\end{enumerate}
\end{proposition}
\lean{New.Data.integral_unfoldStageOn_mul_transposedLiftY, New.Data.integral_mul_liftY_poincareAt, New.Data.sum_quotient_chosenLevel_eq_sum_cutPoints, New.Data.exists_integral_mul_lowerChainUpTo_eq_cutSum, New.Data.exists_integral_eq_mul_inner_restrictToFormY}

\begin{proof}
(a) The descent of $z$ is $[J^\circ_k:J_k]$ times the theta integral of $z$ against the datum of lift $k+1$, because this datum is an arrow for the swapped kernel. Conjugate forms on $Y_{k+1}$ are rapidly decreasing; this does not depend on the twist, and is read off from the trivial twist, whose forms are rapidly decreasing on the source. Hence the second
identity of \cref{prop:B-unfolding}, applied with $W$ the function of stage $k$ at $rg_k$, the datum of lift $k+1$ moved by $g_k$ and $y=z$, gives
\[
  \int_{[R_k]}\mathcal W_k(rg_k)\,(\text{descent }z)(rg_k)\,dr=[J^\circ_k:J_k]\int_{R_{k+1}(E^+)\backslash\U(Y_{k+1})(\Ad)}\mathcal W_{k+1}(h)\,z(h)\,dh ,
\]
both sides converging absolutely; the function $\mathcal W'$ of that proposition is $\mathcal W_{k+1}$, since moving the datum by $g_k$ is the same as evaluating the constant term at $rg_k$.

(b) Lift and descent are transposes, so the left side is $\int_{[\U(Y_0)]}z'\cdot\mathcal P\varphi$ with $z'$ the descent of $z$, a form of level $J_0\subseteq J_M$. By \cref{prop:poincare},
\[
  [J_M:J_0]\int_{[\U(Y_0)]}z'\cdot\mathcal P\varphi=\sum_{\gamma\in J_M/J_0}\int_{[R_0]}\mathcal W_\varphi(r)\,z'(rg\gamma)\,dr ,
\]
and by (a) each term is $[J^\circ_0:J_0]\,\Xi_1(g,\gamma)$; the period depends only on the class of $\gamma$ modulo $J_0$, as $z'$ has level $J_0$. The classes of $J_M/J_0$ are the classes of $\Omega_0$. Since $[J^\circ_0:J_0]/[J_M:J_0]=[J^\circ_0:J_M]$, the claim follows.

(c) Induction on $k$, for all functions $a$ and $S$. For $k=0$, (b) and the hypothesis give $[J^\circ_0:J_M]\,a(\xi)\sum_{c\in\Omega_0}S(c)$, and the multiplicities on $\Omega_0$ are $1$. For the step from $k$ to $k+1$, write the pairing of $z$ with $L_{k+1}\circ\dots\circ L_0(\mathcal P\varphi)$ as the pairing of the descent $z'$ of $z$ with
$L_k\circ\dots\circ L_0(\mathcal P\varphi)$. To apply the induction hypothesis to $z'$, compute the period $\Xi'_{k+1}=\int\mathcal W_{k+1}z'$ by \cref{prop:coset-sum}(b) for the cut $\Omega_{k+1}$ and $f=\mathcal W_{k+1}z'$. The function $f$ is continuous, invariant under the rational points of the Bessel group and integrable by (a); it is supported on
points of bounded orbit by \cref{thm:support}(b), where $m_{k+1}\ne0$; and by (a), with $t$ as the outer variable of index $k+1$, its period over $R_{k+1}$ at $t$ is $[J^\circ_{k+1}:J_{k+1}]$ times $\Xi_{k+2}$ at $(\xi,t_\infty;y,t_f)$, which by hypothesis is $a(\xi,t_\infty)S(y,t_f)$. Therefore
\[
  \Xi'_{k+1}=a'(\xi)\sum_{t\in\Omega_{k+1}}m_{k+1}(t)^{-1}S(y,t),
\]
with $a'(\xi)=\kappa$ the constant of \cref{prop:coset-sum}(b) for the function $t_\infty\mapsto a(\xi,t_\infty)$, which does not depend on $\vartheta$, $\varphi$ or $z$. The induction hypothesis for $z'$, $a'$ and the weighted sum of $S$ over $\Omega_{k+1}$ gives the claim for $k+1$, with the cut sum $\Sigma_{k+2}S$.

(d) The space $Y_m$ is definite and its grading has an empty string, so its Bessel group is all of $\U(Y_m)$ and $F(\iota k)=t(\psi(\mathrm{diag}(k,1)))$ at every point. The space of forms on $Y_m$ is finite-dimensional, the quotient $[\U(Y_m)]$ is a finite union of quotients of the compact real group, and $\bar x$ is the pairing of the algebraic form of $x$ with the
translates of a highest weight vector. Hence the integral is the mean over the finitely many finite classes of an integral over the compact real group of a product of matrix coefficients, which by Schur orthogonality is the pairing divided by the dimension. The map $v\mapsto(k\mapsto t(\psi(\mathrm{diag}(k,1))))$ is equivariant for the block group, so it is a
multiple of the branching map by \cref{thm:B-branching}(c), and the restriction is the branching map divided by a fixed constant. Hence $\int_{[\U(Y_m)]}F\bar x=C\langle x,\mathrm{res}\,\psi\rangle$ with $C$ depending only on $t$.
\end{proof}

The \emph{evaluated tensor} at a tuple $c$ of classes is $\phi_0(y)\otimes\dots\otimes\phi_\varrho(y)$ for representatives $y$ of $c$, and the
\emph{seesaw datum} is the cut sum $\phi_\vartheta=\Sigma_m(\text{evaluated tensor})$, a tensor of finite data for the chain
$X_0\to\dots\to X_m$.
\lean{New.Data.evaluatedTensor, New.Data.evaluatedTensor_mk, New.Data.seesawDatum}

\begin{proposition}\label{prop:seesaw-datum}
\begin{enumerate}[label=\textup{(\alph*)}]
\item $\phi_\vartheta$ is jointly invariant under the levels of $X_0,\dots,X_{m-1}$.
\item There is $c_0$ such that $\iota_\vartheta(\mathcal P\varphi)=c_0\,\mathrm{res}\bigl(\mathrm{chain}(\phi_\vartheta)\varphi\bigr)$ for every twist and every ramified
$\varphi$.
\item There is $C$ such that for every $\vartheta$, $\phi_\vartheta$ is an admissible combination of the $C\,\mathrm{Glue}(x)$ over the admissible
translations, with denominator $(M^{fn_m^2})^{2m}$ away from $q$ and, if $Q=\emptyset$, size $(M^{fn_m^2})^{2\varrho+2}$.
\end{enumerate}
\end{proposition}
\lean{New.Data.isInvariant_seesawDatum, New.Data.exists_lowerChainUpTo_poincareAt_eq, New.Data.exists_isAdmissibleCombination_seesawDatum, New.Data.isAdmissibleCombination_cutSumAt, New.Data.isAdmissibleCombination_cutSum}

\begin{proof}
(a) Let $b$ be a tuple with $b_j$ in the level of $X_j$ for $j<m$ and $b_m=1$, and let $h_j$ be the image of $b_j$ in $\U(Y_j)(\Ad_{\mathrm f})$. By \cref{ss:cuts} the translation by $h_j$ permutes $\Omega_j$ and keeps the multiplicities, so the cut sum of the evaluated tensor at the translated classes $hc$ equals the cut sum at $c$. With
representatives $h_jy_j$ the translated data of lift $j+1$ are moved by $\omega(h_j,h_{j+1})$, and the finite evaluation carries $\omega(\iota b_{2i},\iota b_{2i+1},\iota b_{2i+2})$ on the joined data to the joint action of $(b_{2i},b_{2i+1},b_{2i+2})$ on the two factors of $\phi_i$ (\cref{ss:unfold-pair}). Hence the evaluated tensor at $hc$ is
the joint action of $b$ on the evaluated tensor at $c$, and $\phi_\vartheta$ is fixed by the joint action of $b$.

(b) Pair both sides with a form $x$ on $Y_m$. By \cref{thm:unfolded-upper} for $i=\varrho$, with $y_m=1$, the function $\mathcal W_m$ at $\iota(g)$ is $l(\xi)$ of $(\Theta_{\varrho+1}(y)\varphi)(\iota g)$, where $\Theta_{\varrho+1}(y)\varphi=\mathrm{chain}(\phi(y))\varphi$ for the evaluated tensor $\phi(y)$; by
\cref{prop:outer-integrals}(d) its period against $\bar x$ is $C(l(\xi))\,\langle x,\mathrm{res}\,\mathrm{chain}(\phi(y))\varphi\rangle$, and this depends on $y$ only through its classes, as the evaluated data do not change when each $y_k$ is replaced by a representative of its class. Hence the hypothesis of \cref{prop:outer-integrals}(c) holds with $k+1=m$,
$a(\xi)=C(l(\xi))$ and $S(y)=\langle x,\mathrm{res}\,\mathrm{chain}(\phi(y))\varphi\rangle$, and
\[
  \int_{[\U(Y_m)]}\bar x\cdot\iota_\vartheta(\mathcal P\varphi)=c\,\Sigma_mS=c\,\bigl\langle x,\mathrm{res}\,\mathrm{chain}(\phi_\vartheta)\varphi\bigr\rangle ,
\]
because the cut sum commutes with the linear map $\phi\mapsto\langle x,\mathrm{res}\,\mathrm{chain}(\phi)\varphi\rangle$ (\cref{prop:chains}(b)). The constant $c$ does not depend on $\vartheta$, $\varphi$ or $x$. As the inner product on forms on $Y_m$ is nondegenerate and $x$ is arbitrary, the two forms agree up to the constant $c_0$.

(c) By \cref{thm:good-datum}, the value of the evaluated tensor at a tuple of points integral away from $q$, in particular at representatives of cut classes, is an admissible combination of the $C\,\mathrm{Glue}(x)$ with denominator $(M^{fn_m^2})^{2m}$ and, if $Q=\emptyset$, size $1$. The cut sum is an admissible combination of these values: its
coefficients are products of the weights $m_j(c_j)^{-1}$, and since $m_j$ divides a power of $q$, $1/m_j$ is integral away from $q$, so the denominator is $1$. If $Q=\emptyset$, then $q=1$, the multiplicities are $1$, and each of the $m$ cuts $\Omega_0,\dots,\Omega_{m-1}$ has at most $M^{fn_m^2}$ elements (\cref{ss:cuts}); the cut sum
therefore has size at most $(M^{fn_m^2})^m$, and $m=2\varrho+2$. Substitution multiplies the denominators and the sizes.
\end{proof}

\begin{theorem}[Seesaw identity]\label{thm:seesaw}
Let $t$ be a linear form on the values of forms on $X_0$ and $g$ an archimedean point of $\U(Y_0)$. There is $c\ne0$ such that
every twist $\vartheta$ has a chain datum $\Psi$ with
\[
  \iota_\vartheta(\mathcal P_{g,t}\varphi)=c\,B_{\vartheta,\Psi}(\varphi)\qquad\text{for every ramified source form }\varphi .
\]
\end{theorem}
\lean{New.Data.seesaw_identity}

\begin{proof}
Let $c_0$ and $C$ be as in \cref{prop:seesaw-datum}, $B=M^{fn_m^2}$, and let $u=B^{2m}/(c_0C)$ if $c_0C\ne0$ and $u=1$ otherwise. Write
$\phi_\vartheta=\sum_xr(x)\,C\,\mathrm{Glue}(x)$ as in \cref{prop:seesaw-datum}(c), with $B^{2m}r(x)$ integral away from $q$. If $c_0C\ne0$, then
\[
  u\,c_0\phi_\vartheta=\sum_xB^{2m}r(x)\,\mathrm{Glue}(x)
\]
is an admissible combination of the $\mathrm{Glue}(x)$ with denominator $1$ away from $q$ and, if $Q=\emptyset$, size $B^{2m}B^{2\varrho+2}=M^{3fmn_m^2}$, since $m=2\varrho+2$; if $c_0C=0$, then $u\,c_0\phi_\vartheta=0$, which is the combination with all coefficients $0$. Each admissible translation may be replaced by one supported at the places
dividing $Mq$ with the same glued vacuum: at the other places the translation is integral, as $\bar\pi_M$ is a unit there and the translation is integral away from $q$, and the vacuum is fixed by the integral translations. Collecting the coefficients of equal translations gives a finitely supported function $\Psi$ on tuples of translations. It is
a chain datum (\cref{def:chain-datum}): its support consists of admissible translations supported at the level, its values are integral away from $q$ and bounded by $M^{3fmn_m^2}$ in total if $Q=\emptyset$, and its glued tensor $u\,c_0\phi_\vartheta$ is jointly invariant by \cref{prop:seesaw-datum}(a). By \cref{prop:seesaw-datum}(b) and the
linearity of the chain in its tensor, $u\,\iota_\vartheta(\mathcal P\varphi)=\mathrm{res}(\mathrm{chain}(u\,c_0\phi_\vartheta)\varphi)=B_{\vartheta,\Psi}(\varphi)$, so the theorem holds with $c=u^{-1}$, which is not zero and does not depend on $\vartheta$.
\end{proof}
